\documentclass[11pt]{article}
\usepackage[a4paper,margin=18mm]{geometry}
\usepackage{fontspec}
\setmainfont{Latin Modern Roman}
\usepackage{amsmath,amssymb,amsthm,amscd,graphicx,color}
\usepackage{xurl}
\usepackage[unicode,hidelinks]{hyperref}
\allowdisplaybreaks
\setlength\emergencystretch{3em}

\makeatletter
\newtheorem{theorem}{Theorem}
\newtheorem*{theorem*}{Theorem}
\newtheorem{lemma}{Lemma}
\newtheorem*{lemma*}{Lemma}
\newtheorem{corollary}{Corollary}
\newtheorem*{corollary*}{Corollary}
\newtheorem{proposition}{Proposition}
\newtheorem*{proposition*}{Proposition}
\newtheorem{conjecture}{Conjecture}
\newtheorem*{conjecture*}{Conjecture}
{\theoremstyle{definition} \newtheorem{definition}{Definition}
\newtheorem*{definition*}{Definition}
\newtheorem{example}{Example}
\newtheorem*{example*}{Example}
\newtheorem{remark}{Remark}
\newtheorem*{remark*}{Remark}
\newtheorem{note}{Note}
\newtheorem*{note*}{Note}
}
\makeatother
%\newtheorem{theorem}{Theorem}
\newtheorem{prop}[theorem]{Proposition}
\newtheorem{rem}[theorem]{Remark}
\newtheorem{lem}[theorem]{Lemma}
\newtheorem{def-lem}[theorem]{Definition-Lemma}
\newtheorem{def-prop}[theorem]{Definition-Proposition}
\newtheorem{defin}[theorem]{Definition}
\newtheorem{theor}[theorem]{Theorem}
\newtheorem{cor}[theorem]{Corollary}
\newtheorem{ques}[theorem]{Question}

\def\a{\alpha}\def\Ar{{\rm A}}\def\Ab{\bar{\rm A}}\def\b{\beta}\def\bb{\mathfrak{b}}
%\def\ba{\begin{eqnarray}}\def\ea{\end{eqnarray}}\def\be{\begin{equation}}\def\ee{\end{equation}}
\def\Cprime{C'}\def\ds{\displaystyle}\newcommand{\fracds}[2]{{\ds \frac{#1}{#2}}}
\def\f{\mathfrak{k}}\def\h{\mathfrak{h}}\def\hs{\mathring{z}}\def\g{\mathfrak{g}}\def\sl{\mathfrak{sl}}\def\gl{\mathfrak{gl}}
\def\Ib{\tts\overline{\nns\rm I}_+\tts}\def\J{\mathrm{J}}\def\Jb{\,\overline{\!\rm I}_-\ts}\def\k{\mathfrak{k}}
\def\lb{\label}\def\lcd{\ts,\ldots,}\def\mult{\diamond}\def\n{\mathfrak{n}}\def\nminus{\mathfrak{n}_-}
\def\nplus{\mathfrak{n}}\def\nn{\nonumber}\def\nin{\noindent}
\def\ns{\hskip-1pt}\def\nns{\hskip-.5pt}\def\Mt{M \raisebox{1mm}{$\intercal$}}\def\p{\mathfrak{p}}
\def\pp{\widetilde{\phantom{\,}p\phantom{\,}\phantom{\ }}\phantom{\!}\phantom{\!}\phantom{\!}}
\def\q{\check{\mathrm q}}\def\Q{\mathrm{Q}}\def\r#1{(\ref{#1})}\let\rf=\r
\def\s{z}\def\S{\mathrm {S}}\def\si{\sigma}\def\Sp{\S^{\nplus}}\def\sy{\acute{\sigma}}\def\th{\mathring{h}}
\def\tphi{A}\def\tpsi{A'}\def\tphiplus{B}\def\tpsiminus{B'}\def\ttphi{B}\def\ttpsi{B'}\def\ts{\hskip1pt}\def\tts{\hskip.5pt}
\def\tt{\mathring{t}}\def\U{\mathrm{U}}\def\Uh{\overline{\U}(\h)}\def\UU{\overline{\U}}\def\UUh{{\cal{D}}(\h)}
\def\UUb{\UU^{12}(\bb)}\def\V{\mathrm {V}}\def\ve{\varepsilon}\def\Z{\mathrm{Z}}\def\SZ{\mathrm{SZ}}
\def\Zp{\Z^{\nplus}}\def\z{z}\let\mmult=\mult\newcommand{\CC}{{\mathbb C}}\newcommand{\NN}{{\mathbb N}}
\newcommand{\ZZ}{{\mathbb Z}}\def\Hom{\operatorname{Hom}}\def\I{\mathrm{I}}

%\renewcommand{\theequation}{{\thesection}.{\arabic{equation}}}
%\setcounter{tocdepth}{3}

\numberwithin{equation}{section}


%\newcommand\qedsymbol{\hbox{\rlap{$\sqcap$}$\sqcup$}}\newcommand\qed{\relax\ifmmode\else\unskip\quad\fi\qedsymbol}


\makeatletter\newlabel{seclimit}{{4.4}{}{Original source locator}{}{}}\makeatother
\makeatletter\newlabel{subsection_sl_n}{{4.5}{}{Original source locator}{}{}}\makeatother
\begin{document}
\begin{center}\Large An explicit defining presentation of the diagonal gl n reduction algebra for n at least four over its localized Cartan enveloping algebra\end{center}
\noindent\textbf{Stable ID:} 2573\par
\noindent\textbf{Authors:} Sergei Khoroshkin; Oleg Ogievetsky\par
\noindent\textbf{Original work:} Structure Constants of Diagonal Reduction Algebras of gl Type\par
\noindent\textbf{Version:} Official SIGMA 7 (2011) article 064, DOI 10.3842/SIGMA.2011.064; journal PDF and native TeX acquired 2026-10-01, exact bytes pinned by SHA256\par
\noindent\textbf{Selected task scope:} Theorem3 for n>=\allowbreak{}4: Z\_n is the diagonal gl\_n reduction algebra defined as the double coset of U(gl\_n) tensor U(gl\_n), localized by every h\_ij+\allowbreak{}ell for integer ell and i<j, with extremal-projector multiplication. The six explicitly displayed relation families(4.12),(4.13),(4.14),(4.16),(4.17),(4.18), together with the standing Cartan weight conventions, define the algebra on the weight generators z\_ij,t\_i and are equivalent over the localized Cartan ring to its ordering relations(3.3). Renormalized variables and shifted h are those of Sections2/4.1.\par
\noindent\textbf{Difficulty:} H2: The extremal-projector product defines structure constants by an infinite root-lattice expansion, so ordinary enveloping-algebra PBW relations do not give the required presentation. The authors compute the needed rank-three/four components through the ABRR recurrence, control which weights can survive stabilization, and propagate coefficients by Zhelobenko braid automorphisms. Separate Cartan and opposite-root calculations complete the six relation families before the established presentation criterion can be applied.\par
\noindent\textbf{Editorial boundary note (not author text):} The full original Theorem3 is selected for n at least four with the localized diagonal double-coset algebra, Cartan weight conventions, all shifted/renormalized generators and six displayed relation families. The complete new stabilization, ABRR tensor recurrence, braid induction and all six family derivations are retained before applying the named prior presentation criterion. The author bounded rank-three/four coefficient calculations and stated ancillary reader calculation are preserved. No individual relation, braid helper, cut map or small-rank example is separately counted.\par
\noindent\textbf{Source:} \url{https://sigma-journal.com/2011/064/}\quad DOI: \url{https://doi.org/10.3842/SIGMA.2011.064}\par
\noindent\textbf{License and attribution:} Copyright retained by the named authors. Published by SIGMA under CC BY 4.0: \url{https://creativecommons.org/licenses/by/4.0/}. Source: \url{https://sigma-journal.com/about.html}. Adaptation consists of standalone excerpt formatting, cover and numbering restoration; original author language and mathematical text are retained.\par
\noindent\textbf{Editorial symbol notice (not author text):} The printed symbol U(h) carries the source localization convention; it must not be read as an unlocalized polynomial coefficient ring. The prior presentation criterion[7] is cited background, while all required new relation checks are included in full original pages. No isolated small-rank example is counted.\par
\noindent\textbf{External original locator seclimit:} Section 4.4 discusses a limiting regime for the relation families. This is an introductory pointer to an unselected later interpretation.\par
\noindent\textbf{External original locator subsection\_sl\_n:} Section 4.5 gives an sl\_n variant and its quadratic Casimir; neither is a prerequisite of the selected gl\_n presentation.\par
\medskip\hrule\medskip\noindent\textbf{Author text begins below.}\par
\setcounter{section}{1}
\setcounter{subsection}{0}
\setcounter{subsubsection}{0}
\setcounter{equation}{1}
% Original source lines 114--120
as follows. Suppose $\k$ is given with  a~triangular decomposition
\begin{gather}
\label{intro1} \k=\nminus+\h+\nplus  .
\end{gather}
Denote by $\I_+$ the left ideal of $\Ar :=\U(\g)$, generated by elements of $\nplus$, $\I_+ :=\Ar\nplus$ . Then the reduction algebra $\Sp(\Ar)$, related to the pair
$(\g,\k)$, is def\/ined as the quotient $\mathrm{Norm}(\I_+)/\I_+$ of the normalizer of the ideal $\I_+$ over $\I_+$. It is equipped with a~natural structure of the
associative algebra.  By def\/inition, for any $\g$-module $V$ the space $V^\nplus$ of vectors, annihilated by $\nplus$, is a~module over~$\Sp(\Ar)$. Since $V$ is
\setcounter{section}{1}
\setcounter{subsection}{0}
\setcounter{subsubsection}{0}
\setcounter{equation}{2}
% Original source lines 202--681
\section{Notation}\label{section-notation}

Let ${\mathcal{E}}_{ij}$, $i,j=1\lcd n$, be the standard generators of the Lie algebra $\gl_n$, with the commutation relations
\begin{gather*}
[{\mathcal{E}}_{ij},{\mathcal{E}}_{kl}]=\delta_{jk}{\mathcal{E}}_{il}-\delta_{il}{\mathcal{E}}_{kj} ,
\end{gather*}
where $\delta_{jk}$ is the Kronecker symbol. We shall also use the root notation ${\mathcal{H}}_\alpha$, ${\mathcal{E}}_{\alpha}$, ${\mathcal{E}}_{-\alpha}$, \dots\ for elements of~$\gl_n$.


Let ${\mathcal{E}}_{ij}^{(1)}$ and ${\mathcal{E}}_{ij}^{(2)}$, $i,j=1\lcd n$, be the standard generators of the two copies of the Lie algebra $\gl_n$ in $\g:=\gl_n\oplus\gl_n$,
\begin{gather*}
[{\mathcal{E}}_{ij}^{(a)},{\mathcal{E}}_{kl}^{(b)}]=\delta_{ab}\big(\delta_{jk}{\mathcal{E}}_{il}^{(a)}-\delta_{il}{\mathcal{E}}_{kj}^{(a)}\big) .\end{gather*}
Set
\begin{gather*}
 e_{ij}:={\mathcal{E}}_{ij}^{(1)}+{\mathcal{E}}_{ij}^{(2)} ,\qquad E_{ij}:={\mathcal{E}}_{ij}^{(1)}-{\mathcal{E}}_{ij}^{(2)}
 .
 \end{gather*}
The elements $e_{ij}$ span the diagonally embedded Lie algebra $\k\simeq\gl_n$, while $E_{ij}$ form an adjoint $\k$-module $\mathfrak{p}$. The
Lie algebra $\k$ and the space $\mathfrak{p}$ constitute a symmetric pair, that is, $[\k,\k]\subset \k$, $[\k,\p]\subset \p$, and $[\p,\p]\subset \k$:
\begin{gather*}
 [e_{ij},e_{kl}] = \delta_{jk}e_{il}-\delta_{il}e_{kj}  ,\qquad
  [e_{ij},E_{kl}] = \delta_{jk}E_{il}-\delta_{il}E_{kj}  ,\qquad
   [E_{ij},E_{kl}] = \delta_{jk}e_{il}-\delta_{il}e_{kj} .
\end{gather*}
In the sequel, $h_a$ means the element $e_{aa}$ of the Cartan subalgebra $\h$ of the subalgebra $\k\in\gl_n\oplus\gl_n$ and $h_{ab}$ the
element $e_{aa}-e_{bb}$.


Let $\{\ve_a\}$ be the basis of $\h^*$ dual to the basis $\{ h_a\}$ of $\h$, $\ve_a(h_b)=\delta_{ab}$. We shall use as well the root notation $h_\alpha$,
$e_{\alpha}$, $e_{-\alpha}$ for elements of $\k$, and $H_\alpha$, $E_{\alpha}$, $E_{-\alpha}$ for elements of $\p$.


The Lie subalgebra $\nplus$ in the triangular decomposition \rf{intro1} is spanned by the root vectors~$e_{ij}$ with $i<j$ and the Lie subalgebra $\nminus$ by the
root vectors $e_{ij}$ with $i>j$. Let $\bb_+$ and $\bb_-$
be the corresponding Borel subalgebras, $\bb_+=\h\oplus\nplus$ and $\bb_-=\h\oplus\nminus$. Denote by $\Delta_+$ and $\Delta_-$ the sets of positive and
negative roots in the root system $\Delta=\Delta_+\cup \Delta_-$ of $\k$: $\Delta_+$ consists of roots $\ve_i-\ve_j$ with $i<j$ and $\Delta_-$ consists of roots
$\ve_i-\ve_j$ with $i>j$.
Let $\Q$ be the root lattice, $\Q:=\{\gamma\in\h^*\,|\,\gamma=\sum_{\alpha\in\Delta_+, n_\a\in\ZZ}n_\a \a\}$. It contains the
positive cone $\Q_+$,
\begin{gather*}
\Q_+:=\bigg\{\gamma\in\h^*\,|\,\gamma= \sum_{\alpha\in\Delta_+, n_\a\in\ZZ, \; n_\a\geq 0}n_\a \a\bigg\} .%\label{cqp}
\end{gather*}
For $\lambda,\mu\in\h^*$, the notation
\begin{gather}
\lambda>\mu\label{paor}
\end{gather}
means that the dif\/ference $\lambda-\mu$ belongs to~$\Q_+$, $\lambda-\mu\in\Q_+$. This is a partial order in~$\h^*$.


We f\/ix the following action of the cover of the symmetric group $\S_n$ (the Weyl group of the
diagonal $\k$) on the Lie algebra $\gl_n\oplus\gl_n$ by automorphisms
\begin{gather*}
\sy_i(x):=\mathrm{Ad}_{\exp(e_{i,i+1})}\mathrm{Ad}_{\exp(-e_{i+1,i})}\mathrm{Ad}_{\exp(e_{i,i+1})}(x) ,
\end{gather*}
so that
\begin{gather*}\sy_i(e_{kl})=(-1)^{\delta_{ik}+\delta_{il}}e_{\sigma_i(k)\sigma_i(l)} ,
\qquad\sy_i(E_{kl})=(-1)^{\delta_{ik}+\delta_{il}}E_{\sigma_i(k)\sigma_i(l)}.
\end{gather*}
Here $\sigma_i=(i,i+1)$ is an elementary transposition in the symmetric group. We extend naturally the above action of the cover of $\S_n$ to the action by
automorphisms on the associative algebra $\Ar\equiv\Ar_n:=\U(\gl_n)\otimes \U(\gl_n)$. The restriction of this action to $\h$ coincides
with the natural action $\si(h_{k})=h_{\si(k)}$, $\si\in\S_n$,  of the Weyl group on the Cartan subalgebra.


Besides, we use the shifted action of $\S_n$ on the polynomial algebra $\U(\h)$ (and its localizations) by automorphisms; the shifted action is def\/ined by
\begin{gather}
\label{not1}
\si\circ h_k:= h_{\si(k)}+k-\si(k) ,\qquad k=1,\dots ,n;\qquad \si\in\S_n .
\end{gather}
It becomes the usual action for the variables
\begin{gather}
 \th_k:=h_k-k ,\qquad \th_{ij}:=\th_i-\th_j; \label{hrond}
\end{gather}
by \rf{not1} for any $\si\in\S_n$ we have
\begin{gather*}
\si\circ\th_k=\th_{\si(k)} ,\qquad \si\circ\th_{ij}=\th_{\si(i)\si(j)} .
\end{gather*}

It will be sometimes convenient to denote the commutator $[a,b]$ of two elements $a$ and $b$ of an associative algebra by
\begin{gather}
 \hat{a}(b):=[a,b] .\label{dehaco}
\end{gather}


\section[Reduction algebra $\Z_n$]{Reduction algebra $\boldsymbol{\Z_n}$}\label{section2}

In this section we recall the def\/inition of the reduction algebras, in particular the diagonal reduction algebras of the $\gl$ type.
We introduce the order for which
the ordering relations for the algebra $\Z_n$ will be discussed.
The formulas for the Zhelobenko automorphisms for the algebra $\Z_n$ are given; some basic facts about the standard involution, anti-involution
and central elements for the algebra $\Z_n$ are presented at the end of the section.

{\bf 1.} Let $\Uh$ and $\Ab$ be the rings of fractions of the algebras $\U(\h)$ and $\Ar$ with respect to the multiplicative set, generated by elements
\begin{gather*}
 h_{ij}+l ,\qquad l\in\ZZ  , \quad 1\leq i<j\leq n  .%\label{musel}
 \end{gather*}


Def\/ine $\Z_n$ to be the double coset space of $\Ab$ by its left ideal $\Ib:=\Ab\nplus$, generated by elements of $\nplus$, and the right ideal
$\Jb:=\nminus\Ab$, generated by elements of $\nminus$, $\Z_n:=\Ab/(\Ib+\Jb)$.


The space $\Z_n$ is an associative algebra with respect to the multiplication map
\begin{gather}
\label{not5a}
a\mult b:=aP b  .
\end{gather}
Here $P$ is the extremal projector \cite{AST} for the diagonal $\gl_n$.
It is an element of a certain extension of the algebra $\U(\gl_n)$ satisfying the relations $e_{ij}P=Pe_{ji}=0$ for all $i$ and
$j$ such that $1\leq i<j\leq n$.


The algebra  $\Z_n$ is a particular example of a reduction algebra; in our context, $\Z_n$ is
def\/ined by the coproduct (the diagonal inclusion) $\U(\gl_n)\to\Ar$.

{\bf 2.}  The main structure theorems for the reduction algebras are given in \cite[Section~2]{KO3}.


In the sequel we choose a weight linear basis $\{p_K\}$ of $\p$ ($\p$ is the $\k$-invariant complement to~$\k$ in~$\g$,
$\g =\k +\p$) and equip it with a total order $\prec$. The total order~$\prec$ will be compatible with the partial order $<$ on
$\h^*$, see~\rf{paor}, in the sense that $\mu_K<\mu_L\ \Rightarrow \ p_K\prec p_L$. We shall sometimes write
$I\prec J$ instead of $p_I\prec p_J$. For an arbitrary element $a\in\Ab$ let $\widetilde{a}$ be its image in the reduction
algebra; in particular, $\pp_K$ is the image in the reduction algebra of the basic vector $p_K\in\p$.

{\bf 3.} In our situation we choose the set of vectors $E_{ij}$, $i,j=1,\dots ,n$, as a basis of the space~$\p$. The weight of $E_{ij}$ is $\ve_i-\ve_j$.
The compatibility of a total order $\prec$ with the partial order~$<$ on~$\h^*$ means the condition
\begin{gather*}%\label{not4a}
E_{ij}\prec E_{kl}\qquad \text{if}\quad i-j>k-l .
\end{gather*}
The order in each subset $\{E_{ij}|i-j=a\}$ with a f\/ixed $a$ can be chosen arbitrarily. For instance, we can set
\begin{gather}\label{not4}
E_{ij}\prec E_{kl}\qquad \text{if}\quad i-j>k-l\quad\text{or}\quad i-j=k-l\quad\text{and}\quad i>k .
\end{gather}

Denote the images of the elements $E_{ij}$ in $\Z_n$ by $\s_{ij}$. We use also the notation $t_i$ for the elements $\s_{ii}$ and $t_{ij}:=t_i-t_j$ for the elements
$\s_{ii}-\s_{jj}$. The order \rf{not4} induces as well the order on the generators $\s_{ij}$ of the algebra $\Z_n$:
\begin{gather*} %\label{not4b}
\s_{ij}\prec\s_{kl}\ \Leftrightarrow\ E_{ij}\prec E_{kl}  .
\end{gather*}
The  statement (d) in the paper \cite[Section~2]{KO3} implies an existence of structure constants $\mathrm{B}_{(ab),(cd),(ij),(kl)} \in \Uh$ and $\mathrm{D}_{(ab),(cd)}\in\Uh$
such that for any $a,b,c,d=1,\ldots,n$ we have
\begin{gather}\label{not6}
\s_{ab}\mult\s_{cd}=\sum_{i,j,k,l:\s_{ij}\preceq\s_{kl}}\mathrm{B}_{(ab),(cd),(ij),(kl)}\s_{ij}\mult\s_{kl}+\mathrm{D}_{(ab),(cd)}  .
\end{gather}
In particular, the algebra $\Z_n$ (in general, the reduction algebra related to a symmetric pair $(\k,\p)$, $\g :=\k +\p$) is $\ZZ_2$-graded; the degree of $\s_{ab}$ is~1 and the degree of any element from $\Uh$ is~0.


The relations \rf{not6} together with the weight conditions
\begin{gather*}
 [h,\s_{ab}]=(\ve_a-\ve_b)(h) \s_{ab}%\label{wede}
\end{gather*}
are the def\/ining relations for the algebra $\Z_n$.


Note that the denominators of the structure constants $\mathrm{B}_{(ab),(cd),(ij),(kl)}$ and
 $\mathrm{D}_{(ab),(cd)}$ are pro\-ducts
of linear factors of the form $\th_{ij}+\ell$, $i<j$, where $\ell\geq -1$ is an integer, see~\cite{KO3}.

{\bf 4.} The algebra $\Z_n$ can be equipped with the action of Zhelobenko automorphisms \cite{KO}.  Denote by $\q_{i}$  the Zhelobenko
automorphism $\q_i:\Z_n\to\Z_n$  corresponding to the transposition $\si_i\in\S_n$. It is def\/ined as follows~\cite{KO}. First we def\/ine a map
$\q_i:\Ar\to \Ab/\Ib$ by
\begin{gather} \label{not7}
\q_i(x):=\sum_{k\geq 0}\frac{(-1)^k}{k!} \hat{e}_{i,i+1}^k(\sy_i(x)) e_{i+1,i}^k   \prod_{a=1}^k(h_{i,i+1}-a+1)^{-1}
\quad \mod\Ib  .
\end{gather}
Here $\hat{e}_{i,i+1}$ stands for the adjoint action of the element $e_{i,i+1}$, see~\rf{dehaco}.
The operator $\q_i$ has the property
\begin{gather}
\label{not2}
\q_i(hx)=(\si_i\circ h)\q_i(x)
\end{gather}
for any $x\in\Ar$ and  $h\in\h$; $\si\circ h$ is def\/ined in~\rf{not1}. With the help of~(\ref{not2}), the map $\q_i$ can be extended
to the map (denoted by the same symbol) $\q_i:\Ab\to \Ab/\Jb$ by
setting $\q_i(a(h)x)=(\si_i\circ a(h))\q_i(x)$ for any $x\in\Ar$ and
$a(h)\in\Uh$. One can further prove that $\q_i(\Ib)=0$ and $\q_i(\Jb)\subset (\Jb+\Ib)/\Ib$,
so that $\q_i$ can be viewed as a linear operator $\q_i:\Z_n\to\Z_n$. Due to~\cite{KO},
this is an algebra automorphism, satisfying~\rf{not2}.


The operators $\q_i$ satisfy the braid group relations~\cite{Zh}:
\begin{gather*}
\q_i\q_{i+1}\q_i=\q_{i+1}\q_{i}\q_{i+1}  ,\\
\q_i\q_j=\q_j\q_i  ,\quad |i-j|>1  ,
\end{gather*}
and the inversion relation \cite{KO}:
\begin{gather}
\label{invr}
\q_i^2(x)=\frac{1}{h_{i,i+1}+1}  \sy_i^2(x)  (h_{i,i+1}+1)  ,\qquad x\in\Z_n  .
\end{gather}
In particular, $\q_i^2(x)=x$ if $x$ is of zero weight.

{\bf 5.} The Chevalley anti-involution $\epsilon$ in $\U(\gl_n\oplus\gl_n)$,
$\epsilon(e_{ij}):=e_{ji}$, $\epsilon(E_{ij}):=E_{ji}$, induces the anti-involution $\epsilon$ in
the algebra $\Z_n$:
\begin{gather}
\epsilon(\s_{ij})=\s_{ji}  ,\qquad \epsilon(h_k)=h_k  .\label{anep}
\end{gather}
Besides, the outer automorphism of the Dynkin diagram of $\gl_n$ induces the involutive automorphism
$\omega$ of $\Z_n$,
\begin{gather}
\label{not2a}
\omega(\s_{ij}) =(-1)^{i+j+1}\s_{j'i'}  ,\qquad \omega (h_k)=-h_{k'}  ,
\end{gather}
where $i'=n+1-i$. The operations $\epsilon$ and $\omega$ commute, $\epsilon\omega =\omega\epsilon$.


Central elements of the subalgebra $\U(\gl_n)\otimes 1\subset \Ar$, generated by $n$ Casimir
operators of degrees $1\lcd n$,  as well as central elements of the subalgebra
$1\otimes \U(\gl_n)\subset \Ar$ project to central elements of the algebra $\Z_n$. In particular,
central elements of degree $1$ project to central elements
\begin{gather}
I^{(n,h)}:=h_1+\dots +h_n\label{clih}
\end{gather}
and
\begin{gather}
\label{clit} I^{(n,t)}:=t_1+\dots +t_n
\end{gather}
of the algebra $\Z_n$. The dif\/ference of central elements of degree two projects to the central element
\begin{gather}
\label{drclit}
\sum_{i=1}^n(h_i-2i)t_i
\end{gather}
of the algebra $\Z_n$. The images of other Casimir operators are more complicated.


\section{Main results}

This section contains the principal results of the paper. We f\/irst give
preliminary information on the new basis in which the def\/ining relations for the
algebra $\Z_n$ can be written down in an economical fashion. The braid group action on the new generators is then
explicitly given in Subsection~\ref{brgac}. The complete set of the def\/ining relations for the algebra $\Z_n$ is written down in Subsection~\ref{section3.3}. The regime for which both
the set of the derived def\/ining relations and the set of the def\/ining ordering relation have a controllable ``limiting behavior'' is introduced in Subsection~\ref{seclimit}. Subsection~\ref{subsection_sl_n} deals with the diagonal reduction algebra for $\sl_n$; the quadratic Casimir operator for
${\mathrm{DR}}(\sl_n)$ as well as for the diagonal reduction algebra for an arbitrary semi-simple
Lie algebra~$\k$ is given there. Subsection~\ref{subsection3.4} is devoted to the stabilization and cut phenomena with respect to the embedding of the Lie algebra $\gl_n\oplus\gl_m$ into
the Lie algebra~$\gl_{n+m}$; the theorem about the behavior of the centers of the diagonal reduction algebra under the cutting is proved there.

\subsection{New variables}

We shall use the following elements of $\Uh$:
\begin{gather*}  \tphi_{ij}:=\frac{\th_{ij}}{\th_{ij}-1}  ,\quad \tpsi_{ij}:=\frac{\th_{ij}-1}{\th_{ij}} ,
\quad \ttphi_{ij}:=\frac{\th_{ij}-1}{\th_{ij}-2} ,\quad \ttpsi_{ij}:=\frac{\th_{ij}-2}{\th_{ij}-1} ,\quad \Cprime_{ij}:=\frac{\th_{ij}-3}{\th_{ij}-2},
\end{gather*}
the variables $\th_{ij}$ are def\/ined in~\rf{hrond}. Note that $\tphi_{ij}\tpsi_{ij}=\ttphi_{ij}\ttpsi_{ij}=1$.


Def\/ine  elements $\tt_1,\ldots, \tt_n\in\Z_n$ by
\begin{gather*}
\label{3.1}
\tt_1:=t_1 ,\quad\tt_2:=\q_{1}(t_1) ,\quad\tt_3:=\q_{2}\q_{1}(t_1) ,
\quad\ldots, \quad \tt_n:=\q_{n-1}\cdots \q_{2}\q_{1}(t_1)  .
\end{gather*}
Using \rf{not7} we f\/ind the relations
 \begin{gather}
 \begin{split}
& \q_i(t_i)  =-\frac{1}{\th_{i,i+1}-1}t_i+\frac{\th_{i,i+1}}{\th_{i,i+1}-1}t_{i+1}  ,
\qquad \q_i(t_{i+1})= \frac{\th_{i,i+1}}{\th_{i,i+1}-1}t_i-\frac{1}{\th_{i,i+1}-1}t_{i+1} ,\\ %\nonumber\\
& \q_i(t_k) =  t_k ,\qquad k\not=i,i+1  ,
 \end{split}\label{acwot}
 \end{gather}
which can be used to convert the def\/inition \rf{3.1} into  a linear over the ring $\Uh$ change of variables:
\begin{gather}
\begin{split}
& \ds{\tt_l} =\ds{t_l\prod_{j=1}^{l-1}\tphi_{jl}
-\sum_{k=1}^{l-1}t_k\frac{1}{\th_{kl}-1}\prod_{j=1}^{k-1}\tphi_{jl}}  , \\
& \ds{t_l} =\ds{\tt_l\prod_{j=1}^{l-1}\tpsi_{jl} +\sum_{k=1}^{l-1}\tt_k\frac{1}{\th_{kl}}
\prod_{\substack{j=1 \\ j\neq k}}^{l-1}\tpsi_{jk}} .
\end{split} \label{3.1a}
\end{gather}
For example,
\begin{gather*}
\tt_2 =-\frac{1}{\th_{12}-1}t_1+\frac{\th_{12}}{\th_{12}-1}t_2 ,\qquad
t_2=\frac{1}{\th_{12}}\tt_1+\frac{\th_{12}-1}{\th_{12}}\tt_2 ,\\
\tt_3 =-\frac{1}{\th_{13}-1}t_1-\frac{\th_{13}}{(\th_{13}-1)(\th_{23}-1)}t_2+
\frac{\th_{13}\th_{23}}{(\th_{13}-1)(\th_{23}-1)}t_3  ,\\
t_3 =  \frac{\th_{12}+1}{\th_{12}\th_{13}}\tt_1+\frac{\th_{12}-1}{\th_{12}\th_{23}}\tt_2+
\frac{(\th_{13}-1)(\th_{23}-1)}{\th_{13}\th_{23}}\tt_3  .
\end{gather*}


In terms of the new variables $\tt$'s, the linear in $t$ central element (\ref{clit}) reads
\begin{gather*}
\sum t_i=\sum\tt_i\prod_{a:a\neq i}\frac{\th_{ia}+1}{\th_{ia}} .
\end{gather*}

\subsection{Braid group action}
\label{brgac}

Since
$\q_{i}^2(x)=x$ for any element $x$ of zero weight, the braid group acts as its symmetric group quotient on the space of weight 0 elements.
It follows from \rf{3.1} and $\q_{i}(t_1)=t_1$ for all $i>1$ that
\begin{gather}\label{3.2}
\q_{\sigma}(\tt_i)=\tt_{\sigma(i)}
\end{gather}
for any $\sigma\in \S_n$.


The action of the Zhelobenko automorphisms, see Section \ref{section2},  on the generators $\s_{kl}$  looks as follows:
\begin{alignat}{3}
& \q_i(\s_{ik})=-\s_{i+1,k}\tphi_{i,i+1} , \qquad  &&
\q_i(\s_{ki})=-\s_{k,i+1} ,\quad k\not=i,i+1 ,& \nonumber\\
&\q_i(\s_{i+1,k})=\s_{i,k} ,\qquad && \q_i(\s_{k,i+1})=\s_{k,i}\tphi_{i,i+1} ,\quad k\not=i,i+1 ,& \label{3.5}\\
& \q_i(\s_{i,i+1})=-\s_{i+1,i}\tphi_{i,i+1}\ttphi_{i,i+1} ,\qquad && \q_i(\s_{i+1,i})=-\s_{i,i+1} , & \nonumber\\
&\q_{i}(\s_{j,k})=\s_{j,k} ,\quad j,k\neq i,i+1 .\qquad &&&\nonumber
\end{alignat}

Denote $i'=n+1-i$, as before. The braid group action \rf{3.5} is compatible with the anti-involution $\epsilon$
and the involution $\omega$ (note that $\omega (\th_{ij})=\th_{j'i'}$),  see~\rf{anep} and~\rf{not2a}, in the following sense:
\begin{gather}
\epsilon  \q_i =\q_i^{-1}\epsilon  ,\\
 \omega  \q_i =\q_{i'-1}\omega .\label{qeps}
\end{gather}


Let $w_0$ be the longest element of the Weyl group of $\gl_n$, the symmetric group $\S_n$. Similarly to the squares of
the transformations corresponding to the simple roots, see \rf{invr}, the action of~$\q_{w_0}^2$ is the conjugation by a
certain element of $\Uh$.

\begin{lem}\label{lqw02}
We have
\begin{gather}\label{qw02}
\q_{w_0}^2(x)=S^{-1}xS  ,
\end{gather}
where
\begin{gather}
 S=\prod_{i,j:i<j}\th_{ij}  .\label{conjq02}
 \end{gather}
\end{lem}

The proof shows that the formula \rf{qw02} works for an arbitrary reductive Lie algebra, with $S=\prod_{\alpha\in \Delta_+}\th_{\alpha}$.

\begin{prop}\label{paqw0}
The action of $\q_{w_0}$ on generators reads
\begin{gather}
 \label{3.6}
 \q_{w_0}(\s_{ij})=(-1)^{i+j} \s_{i'j'}\prod_{a:a<i'}\tphi_{ai'}
\prod_{b:b>j'}\tphi_{j'b} ,\\
\q_{w_0}(\tt_i)=\tt_{i'} .\label{3.6a}
\end{gather}
\end{prop}

The proofs of Lemma \ref{lqw02} and Proposition \ref{paqw0} are in Section \ref{sectionproofs}.

\subsection{Def\/ining relations}\label{section3.3}

To save  space we omit in this section the symbol $\mult$ for the multiplication in the algebra~$\Z_n$. It should not lead to any confusion since no other multiplication is used in this section.


Each relation which we will derive will be of a certain weight, equal to a sum of two roots. From general considerations the upper estimate
for the number of terms in a quadratic relation of weight $\lambda=\alpha+\beta$ is the number $|\lambda|$ of quadratic combinations $\s_{\alpha'}\s_{\beta'}$
with $\alpha'+\beta'=\lambda$.
There are several types, excluding the trivial one, $\lambda=2(\ve_i-\ve_j)$, $|\lambda|=1$:
\begin{enumerate}\itemsep=0pt
\item $\lambda=\pm(2\ve_i-\ve_j-\ve_k)$, where $i,j$ and $k$ are pairwise distinct. Then $|\lambda|=2$.
\item $\lambda=\ve_i-\ve_j+\ve_k-\ve_l$ with pairwise distinct $i,j,k$ and $l$. Then $|\lambda|=4$.
\item $\lambda=\ve_i-\ve_j$, $i\neq j$. For $\s_{\alpha'}\s_{\beta'}$, there are $2(n-2)$ possibilities (subtype 3a) with $\alpha'=\ve_i-\ve_k$,
$\beta'=\ve_k-\ve_j$ or $\alpha'=\ve_k-\ve_j$, $\beta'=\ve_i-\ve_k$ with $k\neq i,j$ and $2n$ possibilities (subtype 3b) with $\alpha'=0$,
$\beta'=\ve_i-\ve_j$ or $\alpha'=\ve_i-\ve_j$, $\beta'=0$. Thus $|\lambda|=4(n-1)$.
\item $\lambda=0$. There are $n^2$ possibilities (subtype 4a) with $\alpha'=0$, $\beta'=0$ and $n(n-1)$ possibilities (subtype 4b) with
$\alpha'=\ve_i-\ve_j$,  $\beta'=\ve_j-\ve_i$, $i\neq j$. Here $|\lambda|=n(2n-1)$.
\end{enumerate}

Below we write down relations for each type (and subtype) separately. The relations of the types 1 and 2 have a simple form in terms of the original
generators $\s_{ij}$. To write the relations  of the types 3 and 4, it is convenient to renormalize the generators $\s_{ij}$
with $i\not=j$. Namely, we set
\begin{gather}
\label{not8}\hs_{ij}=\s_{ij}\prod_{k=1}^{i-1}\tphi_{ki}  .
\end{gather}

In terms of the generators $\hs_{ij}$, the formulas \rf{3.5} for the action of the automorphisms $\q_i$  translate as follows:
\begin{alignat*}{3}
& \q_{i}(\hs_{ik})=-\hs_{i+1,k} ,\qquad &&  \q_{i}(\hs_{i+1,k})=\hs_{i,k}\tphi_{i+1,i} ,\quad k\not=i,i+1 , & \\
%\label{not9}
&\q_{i}(\hs_{ki})=-\hs_{k,i+1} ,\qquad && \q_{i}(\hs_{k,i+1})=\hs_{k,i}\tphi_{i,i+1}=\tpsi_{i+1,i}\hs_{k,i} ,\quad k\not=i,i+1,&
\\
&\q_{i}(\hs_{i,i+1})=-\tpsi_{i+1,i}\hs_{i+1,i} , \qquad &&\q_{i}(\hs_{i+1,i})=-\hs_{i,i+1}\tphi_{i+1,i} ,& \\
 &\q_{i}(\hs_{j,k})=\hs_{j,k} , \quad j,k\neq i,i+1 .\qquad &&&
 \end{alignat*}

{\bf 1.} The relations of  the type 1 are:
\begin{gather}\label{relation1}
\s_{ij}\s_{ik}=\s_{ik}\s_{ij}\tphi_{kj} ,\qquad
\s_{ji}\s_{ki}=\s_{ki}\s_{ji}\tpsi_{kj} ,\qquad \text{for}\quad j<k ,\  i\not=j,k .
\end{gather}

{\bf 2.} Denote
\begin{gather*}
D_{ijkl}:=\left( {\ds\frac{1}{\th_{ik}}}-{\ds\frac{1}{\th_{jl}}}\right).
\end{gather*}
Then, for any four pairwise dif\/ferent indices $i$, $j$, $k$ and $l$, we have the following relations of the type 2:
\begin{gather}
\begin{split}
& [ \s_{ij},\s_{kl} ]=\s_{kj}\s_{il}D_{ijkl} , \qquad i<k ,\ j<l  , \\ %\nonumber\\
& \s_{ij}\s_{kl}-\s_{kl}\s_{ij}\tphi_{jl}'\tphi_{lj}'=\s_{kj}\s_{il}
D_{ijkl} ,\qquad i<k ,\ j>l.
\end{split} \label{relation2}
\end{gather}

{\bf 3a.} Let $i\neq k\neq l\neq i$. Denote
\begin{gather*}
\mathring{E}_{ikl}:=-\left((\tt_i-\tt_k) \frac{\th_{il}+1}{\th_{ik}\th_{il}}+(\tt_k-\tt_l) \frac{\th_{il}-1}{\th_{kl}\th_{il}}\right)\hs_{il}+\sum_{a:a\neq i,k,l}
\hs_{al}\hs_{ia} \frac{\ttphi_{ai}}{\th_{ka}+1} .
\end{gather*}
With this notation the f\/irst group of the relations of the type 3 is:
\begin{gather}\notag\hs_{ik}\hs_{kl}\tpsi_{ik} -\hs_{kl}\hs_{ik}\ttphi_{ki}
=\mathring{E}_{ikl} , \qquad i<k<l ,\\   \notag
\hs_{ik}\hs_{kl}\tpsi_{ik}\tpsi_{lk}\ttphi_{lk} -\hs_{kl}\hs_{ik}\ttphi_{ki}
=\mathring{E}_{ikl}  , \qquad i<l<k ,\\
\label{relation3}
\hs_{ik}\hs_{kl}\tphi_{ki} -\hs_{kl}\hs_{ik}\ttphi_{ki}=\mathring{E}_{ikl} , \qquad  k<i<l ,\\
\notag \hs_{ik}\hs_{kl}\tphi_{ki}\tphi_{li}\ttpsi_{li} -\hs_{kl}\hs_{ik}\ttphi_{ki}
=\mathring{E}_{ikl}  ,\qquad  k<l<i ,\\
\notag \hs_{ik}\hs_{kl}\tpsi_{ik}\tpsi_{lk}\ttphi_{lk}\tphi_{li}\ttpsi_{li}-\hs_{kl}\hs_{ik}\ttphi_{ki}
 =\mathring{E}_{ikl}  , \qquad l<i<k , \\  \notag
\hs_{ik}\hs_{kl}\tphi_{ki}\tpsi_{lk}\ttphi_{lk}\tphi_{li}\ttpsi_{li}
-\hs_{kl}\hs_{ik}\ttphi_{ki}=\mathring{E}_{ikl} , \qquad l<k<i .
\end{gather}

The relations \rf{relation3} can be written in a more compact way with the help of both systems, $\s_{ij}$ and $\hs_{ij}$, of generators. Let now
\begin{gather*}
E_{ikl}:=-\left((\tt_i-\tt_k) \frac{\th_{il}+1}{\th_{ik}\th_{il}}+(\tt_k-\tt_l) \frac{\th_{il}-1}{\th_{kl}\th_{il}}\right)
\s_{il}+\sum_{a:a\neq i,k,l}
\hs_{al}\s_{ia} \frac{\ttphi_{ai}}{\th_{ka}+1} .
\end{gather*}
Then
\begin{gather}
\begin{split}
& \s_{ik}\hs_{kl}\tpsi_{ik} -\hs_{kl}\s_{ik}\ttphi_{ki}=E_{ikl} ,\qquad k<l,\\ %\nonumber\\
& \s_{ik}\hs_{kl}\tpsi_{ik}\tpsi_{lk}\ttphi_{lk} -\hs_{kl}\s_{ik}\ttphi_{ki}=E_{ikl} ,\qquad l<k .
\end{split}
\label{shfo3a}
\end{gather}
Moreover, after an extra redef\/inition: $\!\mathring{\hspace{.1cm}z}_{kl}\!\!\!\!\!\!\mathring{\phantom{z}}\ =\hs_{kl}\ttphi_{lk}$ for $k>l$, the left hand side of the
second line in \rf{shfo3a} becomes, up to a common factor, the same as the left hand side of the f\/irst line, namely, it reads $(\s_{ik}
\mathring{\hspace{.1cm}z}_{kl}\!\!\!\!\!\!\mathring{\phantom{z}}\ \ \tpsi_{ik} -\mathring{\hspace{.1cm}z}_{kl}\!\!\!\!\!\!\mathring{\phantom{z}}\ \ \s_{ik}\ttphi_{ki})
\tpsi_{lk}$.

{\bf 3b.} Let $i\neq j\neq k\neq i$. The second group of relations of the type 3 reads:
\begin{gather}\notag\hs_{ij}\tt_{i}=  \tt_i\hs_{ij}\Cprime_{ji}-\tt_j\hs_{ij}  \frac{1}{\th_{ij}+2}-
\sum_{a:a\not=i,j}\hs_{aj}\hs_{ia}  \frac{1}{\th_{ia}+2}  ,\\
\label{3.12}\hs_{ij}\tt_{j}= -\tt_i\hs_{ij}  \frac{\Cprime_{ji}}{\th_{ij}-1}+
\tt_j\hs_{ij}\tphi_{ij}\tpsi_{ji}\ttphi_{ji}+
\sum_{a:a\not=i,j}\hs_{aj}\hs_{ia}\tphi_{ij}\tpsi_{ji}{\ttphi_{ai}}{\th_{ja}+1}  ,\\  \notag
\hs_{ij}\tt_{k}= \tt_i\hs_{ij}  \frac{(\th_{ij}+3)\ttphi_{ji}}{(\th_{ik}^2-1)(\th_{jk}-1)}+
\tt_j\hs_{ij}  \frac{(\th_{ij}+1)\ttphi_{ji}}{(\th_{ik}-1)(\th_{jk}-1)^2}+
\tt_k\hs_{ij}\tphi_{ik}\tphi_{ki}\tphi_{jk}\ttpsi_{jk}\\
  \notag
  \phantom{\hs_{ij}\tt_{k}=}{}
- \hs_{kj}\hs_{ik}  \frac{(\th_{ij}+1)\ttphi_{ki}}{(\th_{ik}-1)(\th_{jk}-1)}
-\sum_{a:a\not=i,j,k}\hs_{aj}\hs_{ia}  \frac{\th_{ij}+1}{(\th_{ik}-1)(\th_{jk}-1)}
%\cdot
\frac{\ttphi_{ai}}{\th_{ka}+1}  .
\end{gather}

{\bf 4a.} The relations of the weight zero (the type~4) are also divided into 2 groups.
This is the f\/irst group of the relations:
\begin{gather}
 \label{relation4a}[\tt_i,\tt_j]=0.
 \end{gather}
As follows from the proof, the relations \rf{relation4a} hold for the diagonal reduction algebra for an arbitrary reductive
Lie algebra: the images of the generators, corresponding to the Cartan subalgebra, commute.

{\bf 4b.} Finally, the second group of the relations of the type 4 is
\begin{gather}
 \label{relation4}
 [\hs_{ij},\hs_{ji}]=\th_{ij}-\frac{1}{\th_{ij}}(\tt_i-\tt_j)^2+
\sum_{a:a\not=i,j}\left(\frac{1}{\th_{ja}+1}\hs_{ai}\hs_{ia}-\frac{1}{\th_{ia}+1}\hs_{aj}\hs_{ja}\right) ,
\end{gather}
where $i\neq j$.

{\bf Main statement.} Denote by $\mathfrak{R}$  the system \rf{relation1},  \rf{relation2}, \rf{relation3},
\rf{3.12},  \rf{relation4a} and \rf{relation4} of the relations.

\begin{theorem} \label{mthe}
The relations $\mathfrak{R}$ are the defining relations for the weight generators $\s_{ij}$ and
$t_i$ of the algebra $\Z_n$. In particular, the set  \eqref{not6} of ordering relations follows over $\Uh$ from $($and is equivalent to$)$  $\mathfrak{R}$.
\end{theorem}

The derivation of the system $\mathfrak{R}$ of the relations is given in Section~\ref{sectionproofs}. The validity in $\Z_n$ of relations from the
set~$\mathfrak{R}$, together with the results
from~\cite{KO3}, completes the proof of Theorem~\ref{mthe} (Section~\ref{proofmain}).

\setcounter{section}{4}
\setcounter{subsection}{5}
\setcounter{subsubsection}{0}
\setcounter{equation}{19}
% Original source lines 789--882
\subsection{Stabilization and cutting} \label{subsection3.4}

In \cite{KO3} we discovered the stabilization and cut phenomena which are heavily used in our derivation of the set of def\/ining relations for the
diagonal reduction algebras of $\gl$-type. The consideration in \cite{KO3} uses the standard (by the f\/irst coordinates) embedding of $\gl_n$ into $\gl_{n+1}$.
In this subsection we shall make several more precise statements about the stabilization and cut considering now the embedding of $\gl_n\oplus\gl_1$
into $\gl_{n+1}$ (more generally, $\gl_n\oplus\gl_m$ into $\gl_{n+m}$). These precisions are needed to establish the behavior of the center of
the diagonal reduction algebra: namely we shall see that cutting preserves the centrality.


Notation: $\h$ in this subsection denotes the Cartan subalgebra of $\gl_{n+m}$.


Consider an embedding of $\gl_n\oplus \gl_m$ into $\gl_{n+m}$, given by an assignment $e_{ij}\mapsto e_{ij}$, $i,j=1,\ldots,n$, and
$e_{ab}\mapsto e_{n+a,n+b}$, $a,b=1,\ldots,m$, where $e_{kl}$ in the source are the generators of $\gl_n\oplus \gl_m$ and target $e_{kl}$
are in $\gl_{n+m}$. This rule together with the similar rule $E_{ij}\mapsto E_{ij}$ and $E_{ab}\mapsto E_{n+a,n+b}$ def\/ines an embedding of the Lie
algebra $(\gl_n\oplus \gl_m)\oplus(\gl_n\oplus \gl_m)$ into the Lie algebra $\gl_{n+m}\oplus\gl_{n+m}$ and of the enveloping algebras
$\Ar_n\otimes \Ar_m=\U(\gl_n\oplus\gl_n)\otimes \U(\gl_m\oplus\gl_m)$ into $\Ar_{n+m}=\U(\gl_{n+m}\oplus\gl_{n+m})$. This embedding clearly
maps nilpotent subalgebras of $\gl_n\oplus \gl_m$ to the corresponding nilpotent subalgebras of $\gl_{n+m}$ and thus def\/ines an embedding
$\iota_{n,m}:\Z_n\otimes\Z_m\to \Z_{n+m}$ of the corresponding double coset spaces. However, the map $\iota_{n,m}$ is not a~homomorphism of
algebras. This is because the multiplication maps are def\/ined with the help of projectors, which are dif\/ferent for $\gl_n\oplus \gl_m$ and $\gl_{n+m}$.


However, as we will explain now we can control certain dif\/ferences between the two multiplication maps. Let $\V_{n,m}$ be the left ideal of the
algebra $\Z_{n+m}$ generated by elements $\s_{ia}$ with $i=1,\ldots,n$ and $a=n+1,\ldots, n+m$; let  $\V_{n,m}'$ be the right ideal of the
algebra $\Z_{n+m}$ generated by elements $\s_{ai}$ with $i=1,\ldots, n$ and $a=n+1,\ldots, n+m$.


Write any element $\lambda\in Q_+$ (the positive cone of the root lattice of $\gl_{n+m}$) in the form $\lambda=\sum_{k=1}^{n+m}\lambda_k\varepsilon_k$.
The element $\lambda$ can be presented as a sum
\begin{gather}\label{lalaplapp}\lambda= \lambda'+\lambda''  ,\end{gather}
where $\lambda'$ is an element of the root lattice of $\gl_n\oplus \gl_m$, and $\lambda''$ is
proportional to the simple root $\varepsilon_n-\varepsilon_{n+1}$: $\lambda'=\sum_{k=1}^{n+m}\lambda'_k\varepsilon_k$ with
$\sum_{k=1}^n\lambda'_k =  \sum_{k=n+1}^{n+m}\lambda'_k=0$ and $\lambda''=c(\varepsilon_n-\varepsilon_{n+1})$.

\begin{lem}\label{lemma1}%\hspace{-.2cm}.
The left ideal  $\V_{n,m}\subset\Z_{n+m}$ consists of images in $\Z_{n+m}$ of  sums
$\sum_{ia}X_{ia} E_{ia}$ with $X_{ia}\in\Ab_{n+m}$, $i=1\lcd n$ and $a=n+1,\dots, n+m$.


The right ideal   $\V_{n,m}\subset\Z_{n+m}$ consists of images in $\Z_{n+m}$ of  sums $\sum_{ai} E_{ai}Y_{ai}$ with $Y_{ai}\in\Ab_{n+m}$,
$i=1\lcd n $ and  $a=n+1
\lcd n+m$.\end{lem}

\begin{proof}
Present the projector $P$ for the Lie algebra $\gl_{n+m}$ as a sum of terms
\begin{gather}
\label{wedeopro} \xi e_{-\gamma_1}\cdots e_{-\gamma_t}e_{\gamma'_{1}}\cdots e_{\gamma'_{t'}}  ,
\end{gather}
where $\xi\in\Uh$, $\gamma_1,\dots ,\gamma_t$ and $\gamma'_1,\dots ,\gamma'_{t'}$ are positive roots of $\gl_{n+m}$. For any $\lambda\in Q_+$
 denote by $P_\lambda$  the sum of above elements with $\gamma_1+\dots +\gamma_{t}=
\gamma'_1+\dots +\gamma'_{t'}=\lambda$. Then $P=\sum_{\lambda\in Q_+}P_\lambda$.
{} For any $X,Y\in\Ab$ def\/ine the element  $X\mult_\lambda Y$ as the image of $XP_\lambda Y$ in the reduction algebra. We have
$X\mult Y=\sum_{\lambda\in Q_+} X\mult_\lambda Y$.


For any $X\in\Ab_{n+m}$, $i=1\lcd n$ and $a=n+1\lcd n+m$ consider the product $X\mult_\lambda \s_{ia}$.


The product $X\mult_\lambda \s_{in}$ is zero if $\lambda''\not=0$ (the component $\lambda''$ is def\/ined by~(\ref{lalaplapp})).
Indeed, in this case in each summand of $P_\lambda$ one of $e_{\gamma'_{k'}}$ is equal to some $e_{jb}$, $j=1\lcd n$ and $b= n+1\lcd n+m$.
Choose an ordered basis of $\n_+$ which ends by all such $e_{jb}$ (ordered arbitrarily);
any element of~$\U(\n_+)$ can be written as a sum of ordered monomials, that is, monomials in which all such~$e_{jb}$ stand on the right.
Since $[e_{jb},E_{ia}]=0$ for any $i,j=1\lcd n$ and $a,b=n+1\lcd n+m$, the product
$e_{\gamma'_{k'}}E_{ia}$ belongs to the left ideal $\Ib$ and thus $X\mult_\lambda \s_{ia}=0$ in $\Z_{n+m}$.


If $\lambda''=0$ then generators of $\n_+$ in monomials entering the decomposition of  $P_\lambda$ are among the elements $e_{ij}$,
$1\leq i<j\leq n$, and $e_{ab}$, $n+1\leq a<b\leq n+m$ and thus their adjoint action leaves the space, spanned by all $E_{ia}$, $i=1\lcd n$,
$a=n+1\lcd n+m$ invariant, so $X\mult_\lambda \s_{ia}$ can be  presented as an image of the sum $\sum_{jb}X_{jb} E_{jb}$ with $X_{jb}\in\Ab_{n+m}$,
$j=1\lcd n$, $b=n+1\lcd n+m$. Thus, the left ideal, generated by all $\s_{ia}$ is contained in the vector space of images in $\Z_{n+m}$ of  sums
$\sum_{jb}X_{jb} E_{jb}$.


Moreover, for any $X\in\Ab_{n+m}$ the element $X\mult \s_{ia}$ is the image of $XE_{ia}+\sum_{j,b:\, j<i,\, b>a}\, X^{(jb)}E_{jb}$ for some $X^{(jb)}$
and the double induction on $i$ and $a$ proves the inverse inclusion.


The second part of lemma is proved similarly.
\end{proof}

\begin{cor}\label{corollary1}
We have the following decomposition of the free left $($and right$)$ $\Uh$-modules:
\begin{gather}\label{st4}
\Z_{n+m}= \mathrm{I}_{n,m}\oplus \Uh\cdot\iota_{n,m}(\Z_n\otimes\Z_m)   ,
\end{gather}
where $\mathrm{I}_{n,m}:=\V_{n,m}+\V'_{n,m}$.
\end{cor}

\begin{proof} The double coset space $\Z_{n+m}$ is a free left and right $\Uh$-module with a basis consisting of images of ordered monomials on elements
$E_{ij}$, $i,j=1\lcd n+m$; recall that we always use orders compatible with the partial order $<$ on $\h^*$, see (c) in Section~\ref{section2},
paragraph~2. We can choose an order for which all ordered monomials are of the form $XYZ$, where $X$ is a monomial on $E_{ai}$ with
$i=1\lcd n$ and $a=n+1\lcd n+m$, $Z$ is a monomial on $E_{ia}$  with $i=1\lcd n$ and $a=n+1\lcd n+m$ while $Y$ is a monomial on
$E_{ij}$ with $i,j=1\lcd n$ or $i,j=n+1\lcd n+m$. Then we apply the lemma above.
\end{proof}
\setcounter{section}{4}
\setcounter{subsection}{6}
\setcounter{subsubsection}{0}
\setcounter{equation}{28}
% Original source lines 1022--1560
\section{Proofs}\label{sectionproofs}


\subsection[Tensor $\J$]{Tensor $\boldsymbol{\J}$}\label{sectionJ}

The multiplication map $\mult$ in $\Z_n$ (we return to the original notation) is given by the
prescription~\rf{not5a}, as in any reduction algebra. It can be formally expanded into
a series over the root lattice of certain bilinear maps as follows. Set
\begin{gather*}
\UU(\bb_\pm):=\Uh\otimes_{\U(\h)}\U(\bb_\pm)  , \qquad \UUb:=\UU(\bb_-)\otimes_{\Uh}  \UU(\bb_+)  .
\end{gather*}
All these are associative algebras. Besides, both algebras~$\UU(\bb_\pm)$ are $\Uh$-bimodules. The
algebra~$\UUb$ admits three commuting actions of~$\Uh$. Two of them are given by the assignments
\begin{gather*}
X(Y\otimes Z):=XY\otimes Z  ,\qquad (Y\otimes Z)X:=Y\otimes ZX  ,
\end{gather*}
for any $X\in\Uh$, $Y\in\UU(\bb_-)$ and $Z\in\UU(\bb_+)$.
The third action associates to any $X\in\Uh$, $Y\otimes Z\in\UUb$ the element $YX\otimes Z=Y\otimes XZ\in\UUb$.


Present the projector $P$ in an ordered form:
\begin{gather} \label{proof1}
P=
\sum_{\gamma,i}\grave{F}_{\gamma,i}\grave{E}_{\gamma,i}\grave{H}_{\gamma,i}=
\sum_{\gamma,i}\grave{H}_{\gamma,i}\grave{F}_{\gamma,i}\grave{E}_{\gamma,i}  ,
\end{gather}
the summation is over $\gamma\in\Q_+$ and $i\in {\mathbb{Z}}_{\geq 0}$;
every $\grave{F}_{\gamma,i}$ is an element of $\U(\n_-)$ of the weight~$-\gamma$, every
$\grave{E}_{\gamma,i}$ is an element of $\U(\n_+)$ of the weight $\gamma$ and
$\grave{H}_{\gamma,i}\in\Uh$. Let $\J$ be the following element of $\UUb$:
\begin{gather*}
\J :=\sum_{\gamma,i}\grave{F}_{\gamma,i}\otimes\grave{E}_{\gamma,i}\grave{H}_{\gamma,i}=
\sum_{\gamma,i}\grave{H}_{\gamma,i}\grave{F}_{\gamma,i}\otimes\grave{E}_{\gamma,i}  ,\qquad \gamma\in\Q_+  ,\quad i\in {\mathbb{Z}}_{\geq 0}  .
\end{gather*}
Due to the PBW theorem in $\U(\gl_n)$ the tensor~$\J$ is uniquely def\/ined by the projector $P$; it is
of total weight zero: $h\J=\J h$ for any $h\in\h$. We have the weight decomposition of $\J$ with
respect to the adjoint action of $\h$ in the second tensor factor of~$\UUb$:
\begin{gather*}
\J= \bigoplus_{\lambda\in\Q_+}\J_\lambda  ,
\end{gather*}
where $\J_\lambda$ consists of all the terms, corresponding to $\grave{F}_{\lambda,i}\grave{E}_{\lambda,i}
\grave{H}_{\lambda,i}$ in~\rf{proof1}
(contributing to $\lambda\in\Q_+$ in the summation),
\begin{gather*}
\J_\lambda :=\sum_i\grave{F}_{\lambda,i}\otimes\grave{E}_{\lambda,i}\grave{H}_{\lambda,i}  .
\end{gather*}
By def\/inition of $\J$, the multiplication $\mult$ in the double coset space $\Z_n$ can be described by the relation
\begin{gather}\label{proof4}
a\mult b= m\left( (a\otimes 1)\J (1\otimes b )\right),
\end{gather}
where $m(\sum_ic_i\otimes d_i)$ is the image in $\Z_n$ of the element $\sum_ic_id_i$.
Moreover, in \rf{proof4} we can replace all products $\grave{E}_{\gamma,i}b$ in the second tensor
factor by the adjoint action of $\grave{E}_{\gamma,i}$ on $b$ (in fact, for $\grave{E}_{\gamma,i}=e_{\gamma_m}\cdots e_{\gamma_1}$,
we can replace $\grave{E}_{\gamma,i}b$ by $[\grave{E}_{\gamma,i},b]$ or by $\hat{e}_{\gamma_m}\cdots \hat{e}_{\gamma_1}(b)$,
see \rf{dehaco}) and likewise all products $a\grave{F}_{\gamma,i}$
in the f\/irst tensor factor by the opposite adjoint action of $\grave{F}_{\gamma,i}$ on $a$.
We have a~decomposition of the product $\mult$ into a sum over $\Q_+$:
\begin{gather}\label{proof5}
a\mult b=\sum_{\lambda\in\Q_+}a\mult_\lambda b  ,\qquad\text{where}\quad
a\mult_\lambda b:=m\left( (a\otimes 1)  \J_\lambda (1\otimes b )\right)  .
\end{gather}
If $a$ and $b$ are weight elements of $\Z_n$ of weights $\nu(a)$ and $\nu(b)$, then the product
$a\mult_\lambda b$ is the image in $\Z_n$ of the sum $\sum_i a_ib_i$, where the weight of each
$b_i$ is $\nu(b)+\lambda$, and the weight of each~$a_i$ is $\nu(a)-\lambda$.


The tensor $\J$ satisf\/ies the Arnaudon--Buf\/fenoir--Ragoucy--Roche (ABRR) dif\/ference equation~\cite{ABRR}, see also~\cite{K} for the translation of the results of~\cite{ABRR} to the language
of reduction algebras. To describe the equation, let  $\vartheta=\frac{1}{2}
\sum_{k=1}^n\th_k^2\in \U(\h)$;
for any positive root $\gamma\in\Delta_+$, denote by~$T_\gamma$ the following linear operator on the vector space~$\UUb$:
\begin{gather*}T_{\gamma}(X\otimes Y):=X e_{-\gamma}\otimes e_{\gamma}Y  .\end{gather*}
The ABRR equation means the relation \cite{ABRR,K}:
\begin{gather*}%\label{proof2}
[1\otimes \vartheta, \J]=-\sum_{\gamma\in\Delta_+}T_{\gamma}(\J)
 .
\end{gather*}
This relation is equivalent to the following system of recurrence relations for the weight components
$\J_\lambda$:
\begin{gather} \label{proof3}
\J_\lambda\cdot\left(\th_\lambda+\frac{(\lambda,\lambda)}{2}\right)=
-\sum_{\gamma\in\Delta_+}T_\gamma\left(\J_{\lambda-\gamma}\right)  ,
\end{gather}
where $\th_\lambda :=\sum_k\lambda_k\th_k$ for $\lambda=\sum_k\lambda_k\ve_k$. The recurrence
relations~(\ref{proof3})  together with the initial condition $\J_0=1\otimes 1$
uniquely determine all weight components~$\J_\lambda$.


It should be noted that the recurrence relations \rf{proof3} provides less information about the structure of the denominators (from~$\U (\h )$)
of the summands of the extremal projector $P$ than the information implied
 by the product formula (see~\cite{AST}) for the extremal projector.


Using \rf{proof3} we get in particular:
\begin{gather} \label{proof3a}
\J_{\a}  = -(\th_{\a}+1)^{-1}e_{-a}\otimes e_{\a} ,
\qquad \a=\ve_i-\ve_{i+1} ,
\\ \notag
 \J_{\a+\beta}  =  (\th_{\a+\b}+1)^{-1}\Bigl(-e_{-\a-\b}\otimes e_{\a+\b}+
(\th_{\a}+1)^{-1} e_{-\a}e_{-\b}\otimes e_{\b}e_{\a}\\
\phantom{\J_{\a+\beta}  =}{}  +(\th_{\b}+1)^{-1} e_{-\b}e_{-\a}\otimes e_{\a}e_{\b}\Bigr)
  ,\qquad \a=\ve_{i-1}-\ve_{i}  ,\quad \b=\ve_{i}-\ve_{i+1}  ,\label{proof3b}
\\
  \label{proof3c}\J_{\ve_i-\ve_{j}+\ve_k-\ve_{l}}= \J_{\ve_i-\ve_{j}}\cdot
\J_{\ve_k-\ve_{l}}  ,\qquad  i<j<k<l  .
\end{gather}

\subsection{Braid group action}

%{\bf 1.}
The proof of the relations \rf{acwot} and \rf{3.5} consists of the following arguments, valid for any
reduction algebra. Let $\a$ be any simple root of $\gl_n$, $\a=\ve_i-\ve_{i+1}$ and $\g_\a$ the
corresponding $\sl_2$ subalgebra of $\gl_n$. It is spanned by the elements $e_\a=e_{i,i+1}$,
$e_{-\a}= e_{i+1,i}$ and $h_\a= h_i-h_{i+1}$. Let $\sy_\a=\sy_i$ be the corresponding
automorphism of the algebra $\Ar$ and $\q_\a=\q_i$ the Zhelobenko automorphism of
$\Z_n$. Assume that $Y\in\Ar$ belongs, with respect to the adjoint action of $\g_\a$, to an irreducible
f\/inite-dimensional $\g_\a$-module of dimension $2j+1$, $j\in\{\ts0,1/2,1,\ldots\ts\}$. Assume further
that $Y$ is homogeneous, of weight $2\ts m$, $[h_\a,Y]=2m Y$. Identify $Y$ with its image in $\Z_n$. Then
$\q_\a(Y)$ coincides with the image in $\Z_n$ of the element
\begin{gather*}%\label{proof6}
\prod_{i=1}^{j+m}\,(h_\a+i+1)\ts\cdot\ts\sy_\alpha(Y)\cdot
{\ds\prod_{i=1}^{j+m} (h_\a-i+1)^{-1}}  .
\end{gather*}
This can be checked directly using \cite[Proposition 6.5]{KO}.


In the realization of irreducible $\sl_2$-modules as the spaces of homogeneous polynomials in two
variables $u$ and $v$,
\begin{gather*}e_\a\mapsto u\frac{\partial}{\partial v}  , \qquad h_\a\mapsto u\frac{\partial}{\partial u}-v\frac{\partial}{\partial v}
\qquad {\mathrm{and}}\qquad  e_{-\a}\mapsto v\frac{\partial}{\partial u}  ,
\end{gather*}
the operator $\sy_\alpha$ becomes $(\sy_\alpha f)(u,v)=f(-v,u)$, or, in the basis $|j,k\rangle :=x^{j+k}y^{j-k}$ ($j$ labels the
representation; $k=0,1,\dots ,2j$),
\begin{gather*}\sy_\alpha : \  |j,-j+k\rangle\mapsto (-1)^k |j,j-k\rangle  .\end{gather*}

\begin{proof}[Proof  of Lemma \ref{lqw02}, Subsection~\ref{brgac}.] %, Subsection \ref{brgac}.
To see this, write a reduced expression for $\q_{w_0}$, $\q_{w_0}=\q_{\alpha_{i_1}}\cdots \q_{\alpha_{i_M}}$
with $\alpha_{i_1},\dots ,\alpha_{i_M}$ simple roots. Then $\q_{w_0}=\q_{\alpha_{i_M}}\cdots \q_{\alpha_{i_1}}$ as well.
Writing, for $\q_{w_0}^2$, the second expression after the f\/irst one, we get squares of
$\q_{\alpha_{i_s}}$'s (which are conjugations by $\th_{\alpha_{i_s}}^{-1}$'s; they thus commute) one after another.
Moving these conjugations to the left through the remaining~$\q$'s, we produce, exactly like in the
construction of a system of all positive roots from a reduced expression for
the longest element of the Weyl group of a reductive Lie group, the conjugation by the product (\ref{conjq02})
over all positive roots.
\end{proof}

\begin{proof}[Proof  of Proposition~\ref{paqw0}, Subsection~\ref{brgac}.] %, subsection \ref{brgac}.
Only formula \rf{3.6} needs a proof (formula~\rf{3.6a}
is a particular case of~\rf{3.2}).


For a moment, denote the longest element
of the symmetric group $S_n$ by $\q_{w_0}^{(n)}$. Let $\psi_j:=\q_j\q_{j-1}\cdots\q_1$ (the product in the descending order).
We have $\q_{w_0}^{(n+1)}=\q_{w_0}^{(n)}\psi_n$ and  $\q_{w_0}^{(n+1)}=\psi_1\psi_2\cdots\psi_n$ (the product in the
ascending order).


For $j<n$ it follows from \rf{3.5} that $\psi_j (\s_{n+1,1})=(-1)^j\s_{n+1,j+1}$ (say, by induction on $j$). So,
\begin{gather*}\psi_n(\s_{n+1,1})=q_n\psi_{n-1} (\s_{n+1,j+1})=(-1)^{n-1}q_n(\s_{n+1,n})=(-1)^{n}\s_{n,n+1}  ,\end{gather*}
again by \rf{3.5}. Next, $\psi_k\psi_{k+1}\cdots\psi_{n-1}(\s_{n,n+1})=\s_{k,n+1}$ by induction on $n-k$ and
again \rf{3.5}. Thus,
\begin{gather}\label{binh}
\q_{w_0}(\s_{n+1,1})=(-1)^n \s_{1,n+1} ,
\end{gather}
establishing \rf{3.6} for $i=n+1$ and $j=1$. We now prove \rf{3.6} for $i>j$ (positions below the main diagonal) by induction backwards
on the height $i-j$ of a negative root;  the formula~\rf{binh} serves as the induction base. Assume that~\rf{3.6}
is verif\/ied for a given level $i-j$ and $i-j-1>0$ (so that the positions $(i,j+1)$ and $(i-1,j)$ are still under the main
diagonal). By \rf{3.5}, $\s_{i,j+1}=-\q_j(\s_{ij})$, therefore
\begin{gather*}\q_{w_0}(\s_{i,j+1}) = -\q_{w_0}(\q_j(\s_{ij}))=-\q_{j'-1}(\q_{w_0}(\s_{ij})) \\
\phantom{\q_{w_0}(\s_{i,j+1})}{}  = (-1)^{i+j+1} \q_{j'-1}\left(\s_{i'j'}\ds{\prod_{a:a<i'}}\tphi_{ai'}  \ds{\prod_{b:b>j'}}\tphi_{j'b}\right)\\
\phantom{\q_{w_0}(\s_{i,j+1})}{} = (-1)^{i+j+1}\s_{i',j'-1}\tphi_{j'-1,j'}\ds{\prod_{a:a<i'}}\tphi_{ai'}  \ds{\prod_{b:b>j'}}\tphi_{j'-1,b}\\
\phantom{\q_{w_0}(\s_{i,j+1})}{} = (-1)^{i+j+1}\s_{i',(j+1)'}\ds{\prod_{a:a<i'}}\tphi_{ai'}\ds{\prod_{b:b>(j+1)'}}\tphi_{(j+1)',b}  .
\end{gather*}
In the second equality we used the identity $\q_{w_0}\q_j=\q_{j'-1}\q_{w_0}$ in the braid group; the third equality is the
induction assumption; in the fourth equality we used that $i'\neq j'-1$ (since $i-j-1>0$) and then~\rf{3.5}; in the f\/ifth equality we
replaced $j'-1$ by $(j+1)'$. The calculation for $\q_{w_0}(\s_{i-1,j})$ is similar; it uses $\s_{i-1,j}=\q_{i-1}(\s_{ij})$.
The proof of the formula~\rf{3.6} for positions below the main diagonal is f\/inished.


The proof of  \rf{3.6} for $i<j$ (positions above the main diagonal) follows now from Lem\-ma~\ref{lqw02}. %, Subsection \ref{brgac}.
\end{proof}

\subsection{Derivation of relations}\label{sectionDerRel}

The set  of def\/ining relations in $\Z_n$ divides into several dif\/ferent types,
see Section~\ref{section3.3}. We prove the necessary amount of relations
of each type and get the rest by applying the transformations from the braid group as well as the anti-involution~$\epsilon$, see~(\ref{anep}).


We never use the automorphism $\omega$, def\/ined in \rf{not2a}, in the derivation of relations. However, the involution $\omega$ is compatible with our set
of relations in the sense explained in Section \ref{proofmain}.


In the following we denote by the symbol $ \equiv $ the equalities of elements from $\Ab$ modulo the sum $(\Jb+\Ib )$ of two ideals $\Jb$ and $\Ib$
def\/ined in the beginning of Section \ref{section2}. Moreover,
for any two elements $X$ and $Y$ of the algebra $\Ab$ we may regard the expressions $X\mult Y$ and $X\mult_\lambda Y$ as the sums
of elements from $\Ab$ def\/ined in \rf{proof4} and \rf{proof5}. The sum $X\mult_\lambda Y$  is f\/inite. By the construction,
all but a f\/inite number of terms in the product $X\mult Y$ belong to $(\Jb+\Ib )$. Unlike to the system of notation
adopted in Section~\ref{section2}, our proof of each relation in $\Z_n$ will use equalities in $\Ab$ taken modulo $(\Jb+\Ib )$.


We also use the notation $H_i$ for the element $E_{ii}\in\Ar$ and $H_{ij}=H_i-H_j=E_{ii}-E_{jj}$.

{\bf 1.} We f\/irst prove in $\Z_n$ the relation
\begin{gather}\label{proof7}
\s_{12}\mult\s_{13}=\s_{13}\mult\s_{12}  \frac{\th_{23}}{\th_{23}+1}  .
\end{gather}
Elements $\s_{12}$ and $\s_{13}$ are images in $\Z_n$ of $E_{12}$ and $E_{13}$. Consider the
product $E_{12}\mult_\lambda E_{13}$. Since the adjoint action of $\gl_n$ preserves the space
$\p$, see Section~\ref{section-notation}, this product is the sum of such monomials $E_{ij}E_{kl}$,
with coef\/f\/icients in $\Uh$, that (i): the weight $\ve_k-\ve_l$
of $E_{kl}$ is equal to the weight $\ve_{1}-\ve_3$ of $E_{13}$ plus $\lambda\in\Q_+$, while (ii): the
weight $\ve_i-\ve_j$ of $E_{ij}$ is equal to the weight $\ve_{1}-\ve_2$ of $E_{12}$ minus
$\lambda$. Assume that $E_{12}\mult_\lambda E_{13}\neq 0$. By (i), $\lambda =-\ve_1+\ve_3+\ve_k-\ve_l$
and it can be positive only if $k=1$ and $l\geq 3$. So, the condition (i) implies that either
$\lambda=0$ or $\lambda=\ve_3-\ve_l$ with $l>3$. The possibility $\lambda=\ve_3-\ve_l$, $l>3$, is excluded
by the condition (ii). Therefore, $\lambda=0$ and
\begin{gather}
\label{proof8}
E_{12}\mult E_{13}\equiv E_{12}E_{13} .
\end{gather}
Similarly, for $\lambda\in\Q_+$, which can non-trivially contribute to the product $E_{13}\mult E_{12}$,
the analogue of the condition (i) on the weight $\lambda$ gives the
restriction $\lambda =0$ or $\lambda=\ve_2-\ve_k$, $k>2$; the analogue of the condition (ii)
further restricts $\lambda$:  $\lambda=0$ or $\lambda=\ve_2-\ve_3$, so we have
\begin{gather*}
E_{13}\mult_{\ve_2-\ve_3} E_{12}\equiv -E_{13}e_{32}e_{23}  \frac{1}{\th_{23}+1}E_{12}\equiv
-E_{13}e_{32}e_{23}E_{12}  \frac{1}{\th_{23}}\equiv E_{12}E_{13}  \frac{1}{\th_{23}}  ,
\end{gather*}
since $\J_{\ve_2-\ve_3}=-e_{32}\otimes e_{23}(\th_{23}+1)^{-1}$ as it follows from the ABRR
equation, see~(\ref{proof3a}), or from the precise explicit expression for the projector~$P$,
see~\cite{AST}. Thus, since $E_{12}$ and $E_{13}$ commute in the universal enveloping algebra
\begin{gather}\label{proof9}E_{13}\mult E_{12}\equiv E_{13}E_{12}+E_{13}\mult_{\ve_2-\ve_3} E_{12}=
E_{12}E_{13}\left(1+\frac{1}{\th_{23}}\right),
\end{gather}
Comparing \rf{proof8} and \rf{proof9} we f\/ind \rf{proof7}.


Applying to \rf{proof7} the anti-involution $\epsilon$, see (\ref{anep}), we get the relation
\begin{gather}\label{proof10}
\s_{21}\mult\s_{31}=\s_{31}\mult\s_{21}  \frac{\th_{23}+1}{\th_{23}}  .
\end{gather}

The rest of the relations \rf{relation1} are obtained from~\rf{proof7} and~\rf{proof10}
by applying dif\/ferent transformations $\q_w$ from the Weyl group.

{\bf 2.} Now we prove in $\Z_n$ the relation
\begin{gather}\label{proof11}
\s_{13}\mult\s_{24}-\s_{24}\mult\s_{13}=
\left(\frac{1}{\th_{12}}-\frac{1}{\th_{34}}\right)\s_{23}\mult\s_{14}  .
\end{gather}
We begin by the proof of this relation in $\Z_4$. We proceed in the same manner as for the derivation of the
relation (\ref{proof7}),
\begin{gather*}E_{13}\mult E_{24}\equiv E_{13}E_{24}+E_{13}\mult_{\ve_1-\ve_2} E_{24} \\
\phantom{E_{13}\mult E_{24}}{} \equiv
E_{13}E_{24}-E_{13}e_{21}e_{12}  \frac{1}{\th_{12}+1} E_{24}
 \equiv E_{13}E_{24}+E_{23}E_{14} \frac{1}{\th_{12}}  ,\\
 E_{24}\mult E_{13}\equiv E_{24}E_{13}+E_{24}\mult_{\ve_3-\ve_4} E_{13}\\
 \phantom{E_{24}\mult E_{13}}{}  \equiv
E_{24}E_{13}-E_{24}e_{43}e_{34}  \frac{1}{\th_{34}+1} E_{13} \equiv E_{13}E_{24}+E_{23}E_{14}  \frac{1}{\th_{34}}  ,\\
 E_{23}\mult E_{14} \equiv E_{23}E_{14}  .
 \end{gather*}
Combining the three latter equalities we obtain \rf{proof11} in $\Z_4$.


The dif\/ference of the left and right hand sides of \rf{proof11} in $\Z_n$ is a linear combination of monomials in $\s_{ij}$ of
the total weight $\ve_1+\ve_2-\ve_3-\ve_4$. The weight is non-trivial, so the monomials can be only quadratic.
Due to the stabilization phenomenon, each monomial should contain $\s_{ij}$ with $i>4$ or $j>4$, but, by the weight arguments,
there is no such non-zero possibility, which completes the proof of the relation \rf{proof11} in $\Z_n$.


The rest of relations \rf{relation2} is then obtained by applications of the transformations from the braid group.

{\bf 3a.} We continue and derive in $\Z_4$ the relation (we remind the notation $t_{ij}:=\s_{ii}-\s_{jj}$,
see Section~\ref{section2},
and $H_{ij}=E_{ii}-E_{jj}$):
\begin{gather} \label{proof12}
\s_{23}\mult\s_{12}-\s_{12}\mult\s_{23}=t_{12}\mult\s_{13}
 \frac{1}{\th_{12}}+t_{23}\mult\s_{13}  \frac{1}{\th_{23}}-\s_{43}\mult\s_{14}  \frac{\th_{34}+1}{\th_{34}\th_{24}}  .
\end{gather}
Using \rf{proof3a}--\rf{proof3c}, we calculate, to obtain the result for $\Z_4$:
\begin{gather}E_{12}\mult E_{23} \equiv E_{12}E_{23}+ E_{12}\mult_{\ve_1-\ve_2} E_{23}+
E_{12}\mult_{\ve_1-\ve_2+\ve_3-\ve_4} E_{23} \nonumber\\
\phantom{E_{12}\mult E_{23}}{}
 \equiv E_{12}E_{23}-H_{12}E_{13} \frac{1}{\th_{12}}  ,\label{1223r1}\\
  E_{23}\mult E_{12} \equiv E_{23}E_{12}+ E_{23}\mult_{\ve_2-\ve_3} E_{12}+
E_{23}\mult_{\ve_2-\ve_4} E_{12}\notag\\
\phantom{E_{23}\mult E_{12}}{} \equiv E_{23}E_{12}+H_{23}E_{13}  \frac{1}{\th_{23}}-E_{43}E_{14} \frac{(\th_{23}-1)}{\th_{23}
\th_{24}} ,\label{1223r2}\\
  H_{12}\mult E_{13} \equiv H_{12} E_{13}+H_{12}\mult_{\ve_3-\ve_4} E_{13}\equiv
H_{12} E_{13}  ,\label{1223r3}\\
  H_{23}\mult E_{13} \equiv H_{23} E_{13}+H_{23}\mult_{\ve_3-\ve_4} E_{13}\equiv
H_{23} E_{13}+E_{43}E_{14}  \frac{1}{\th_{34}}  ,\label{1223r4}\\
  E_{43}\mult E_{14} \equiv E_{43}E_{14} .\label{1223r5}
  \end{gather}
Combining the above equalities and taking into account that $[E_{12},E_{23}]=e_{13}\equiv 0$,
we get~\rf{proof12} in~$\Z_4$. We could apply here the stability arguments (as we shall do in the sequel)
but we give some more details at this point to give a f\/lavor of how such derivations of relations work.
For the same, as \rf{1223r1}--\rf{1223r5}, calculations for $\Z_n$, the analogues of the conditions (i)
and (ii), see paragraph 1 of this subsection, restrict $\lambda$ to be of the form $\ve_1-\ve_2+\ve_3-\ve_k$, $k\geq 3$ for
\rf{1223r1}; $\ve_2-\ve_k$, $k\geq 2$ for \rf{1223r2}; $\ve_3-\ve_k$, $k\geq 3$ for~\rf{1223r3}
and \rf{1223r4}; $\ve_4-\ve_k$, $k\geq 4$ for~\rf{1223r5}. It follows, for, say, $n=5$, that the right
hand sides of \rf{1223r1}--\rf{1223r5} might be modif\/ied only by an addition of the term proportional to
$E_{53}E_{15}$; and this will be compensated by an addition of the term, proportional to
$\s_{53}\mult\s_{15}$ to the right hand side of~\rf{proof12}, since $E_{53}\mult E_{15}\equiv E_{53}E_{15}$;
the proportionality coef\/f\/icient is uniquely def\/ined.
This pattern clearly continues and we conclude that there is a relation in $\Z_n$ of the form
\begin{gather}\label{proof13}
\s_{23}\mult\s_{12}-\s_{12}\mult\s_{23}=t_{12}\mult\s_{13} \frac{1}{\th_{12}}+t_{23}\mult\s_{13}\frac{1}{\th_{23}}-
\sum_{k>3}\s_{k3}\mult\s_{1k}X_k  ,
\end{gather}
with certain, uniquely def\/ined, coef\/f\/icients $X_k\in\Uh$, $k=4\lcd n$, and already known $X_4=(\th_{34}+1)
{\th_{34}^{-1}\th_{24}^{-1}}$. To f\/ind $X_5,\dots ,X_n$, we apply to \rf{proof13} the
automorphisms $\q_k$, $k=4\lcd n-1$, which leave invariant the left hand side and the f\/irst two terms in the
right hand side of~\rf{proof13}. The uniqueness of the relation of the form~\rf{proof13}, together with the equality
$\q_k(\s_{k3}\mult\s_{1k})=\s_{k+1,3}\mult\s_{1,k+1}(\th_{k,k+1}+1)\th_{k,k+1}^{-1}$, imply the recurrence relation
$X_{k+1}=\q_k(X_k)\cdot(\th_{k,k+1}+1)\th_{k,k+1}^{-1}$ and we f\/ind
\begin{gather*}%\label{proof14}
X_k=\frac{1}{\th_{2k}}\prod_{j=3}^{k-1}\frac{\th_{jk}+1}{\th_{jk}}  .
\end{gather*}

After the renormalization~\rf{not8} and the change of variables~\rf{3.1a}, the derived relation becomes one of the
relations in the f\/irst line of~\rf{relation3}.


Applying the transformations from the braid group, we obtain the rest of the relations from the list~\rf{relation3}.

{\bf 3b.} We have the following equalities in $\Z_3$:
\begin{gather}\label{proof15}
\s_{12}\mult t_1=t_1\mult\s_{12}  {\ds\frac{\th_{12}+2}{\th_{12}+1}}-t_2\mult\s_{12}  {\ds\frac{1}{\th_{12}+1}}
-\s_{32}\mult\s_{13}  {\ds\frac{\th_{23}+1}{\th_{23}(\th_{13}+1)}}  ,\\
\label{proof16}
\s_{12}\mult t_2=-t_1\s_{12}  {\ds\frac{1}{\th_{12}+1}}+t_2\mult\s_{12}
 {\ds\frac{\th_{12}+2}{\th_{12}+1}}+\s_{32}\mult\s_{13}  {\ds\frac{\th_{13}+2}{\th_{23}(\th_{13}+1)}}  ,
\end{gather}
and the equality in $\Z_4$:
\begin{gather}\label{proof17}
\s_{12}\mult t_4=t_4\mult\s_{12}-\s_{42}\mult\s_{14}  {\ds\frac{\th_{12}+1}{(\th_{14}+1)\th_{24}}}  .
\end{gather}
Equalities \rf{proof15} and \rf{proof16} are the results of the following calculations for $\Z_3$,
using \rf{proof3a}--\rf{proof3c}, and of the commutativity $[H_1,E_{12}]=e_{12}\equiv 0$,
$[H_2,E_{12}]=-e_{12}\equiv 0$:
\begin{gather*}E_{12}\mult H_1 \equiv E_{12}H_1+H_{12}E_{12}  \frac{1}{\th_{12}+1}-
E_{32}E_{13}  \frac{\th_{12}}{(\th_{12}+1)(\th_{13}+1)}  , \\
 E_{12}\mult H_2 \equiv E_{12}H_2-H_{12}E_{12}  \frac{1}{\th_{12}+1}-
E_{32}E_{13}  \frac{1}{(\th_{12}+1)(\th_{13}+1)}, \\
 H_1\mult E_{12} \equiv H_1E_{12}  , \qquad
 H_2\mult E_{12} \equiv H_2E_{12}-E_{32}E_{13}  \frac{1}{\th_{23}} , \qquad
 E_{32}\mult E_{13} \equiv E_{32}E_{13}  .\end{gather*}
The derivation of \rf{proof17} can be done with the help of the following calculations for $\Z_4$:
\begin{gather}\label{pr14n}
E_{12}\mult H_4 \equiv E_{12}H_4+E_{42}E_{14}  \frac{1}{\th_{14}+1},\\
 H_4\mult E_{12} \equiv H_4E_{12}+E_{42}E_{14}  \frac{1}{\th_{24}}  ,\qquad
 E_{42}\mult E_{14} \equiv E_{42}E_{14}  .\nonumber
 \end{gather}
We shall make a comment about the line \rf{pr14n} only. Here one might expect, by  the analogues of the conditions (i) and (ii), see paragraph 1 of this subsection,
non-trivial contributions to $E_{12}\mult H_4$ from the weights 0, $\ve_1-\ve_2$, $\ve_1-\ve_3$ and $\ve_1-\ve_4$.
So we need, in addition to \rf{proof3a}--\rf{proof3c}, some information about
$\J_{\ve_1-\ve_4}$. It follows from the ABRR equation that $\J_{\ve_1-\ve_4}(\th_{14}+1)=-T_{\ve_1-\ve_2}(\J_{\ve_2-\ve_4})
-T_{\ve_1-\ve_3}(\J_{\ve_3-\ve_4})-T_{\ve_1-\ve_4}(\J_0)-T_{\ve_2-\ve_3}(\J_{\ve_1-\ve_2+\ve_3-\ve_4})-
T_{\ve_2-\ve_4}(\J_{\ve_1-\ve_2})-T_{\ve_3-\ve_4}(\J_{\ve_1-\ve_3})$. Since $e_{13}$ and $e_{12}$
commute with $H_4$, the parts $T_{\ve_2-\ve_4}(\J_{\ve_1-\ve_2})$ and $T_{\ve_3-\ve_4}(\J_{\ve_1-\ve_3})$
do not contribute; $e_{42}$ and $e_{43}$ commute with $E_{12}$, so the parts  $T_{\ve_1-\ve_2}(\J_{\ve_2-\ve_4})$
and $T_{\ve_1-\ve_3}(\J_{\ve_3-\ve_4})$ do not contribute either; $\J_{\ve_1-\ve_2+\ve_3-\ve_4}=
\J_{\ve_3-\ve_4}\J_{\ve_1-\ve_2}$ does not contribute again since $e_{12}$ commute with $H_4$.
Thus the only contribution is from $T_{\ve_1-\ve_4}(\J_0)$ and we quickly arrive at~\rf{pr14n}.


Applying the automorphism $\q_3$ of the algebra $\Z_4$ to the relation \rf{proof17}, see \rf{acwot} and~\rf{3.5}, we f\/ind
\begin{gather*}
\s_{12}\mult\big(t_3\th_{34}-t_4\big)=\big(t_3\th_{34}-t_4\big)\mult \s_{12}
-\s_{32}\mult\s_{13}  {\ds\frac{\th_{34}(\th_{12}+1)}{(\th_{13}+1)\th_{23}}}  .
\end{gather*}
We then add \rf{proof17} to this relation  and obtain the following relation in $\Z_4$:
\begin{gather}\label{proof18}
\s_{12}\mult t_3=t_3\mult\s_{12}
-\s_{32}\mult \s_{13}  {\ds\frac{(\th_{12}+1)}{(\th_{13}+1)\th_{23}}}-
\s_{42}\mult\s_{14}  {\ds\frac{\th_{12}+1}{(\th_{14}+1)\th_{24}\th_{34}}}  .
\end{gather}

The stabilization arguments for \rf{proof15}, \rf{proof16} and \rf{proof18} imply the existence of the following relations in $\Z_n$:
\begin{gather}\label{proof19}\s_{12}\mult t_1 =t_1\mult\s_{12}  {\ds\frac{\th_{12}+2}{\th_{12}+1}}-
t_2\mult\s_{12}  {\ds\frac{1}{\th_{12}+1}}+\sum_{k>2}\s_{k2}\mult\s_{1k}X_k^{(1)},\\
 \label{proof20}\s_{12}\mult t_2 =-t_1\mult\s_{12}  {\ds\frac{1}{\th_{12}+1}}+t_2\mult\s_{12}  {\ds\frac{\th_{12}+2}{\th_{12}+1}}
+\sum_{k>2}\s_{k2}\mult\s_{1k}X_k^{(2)},
\\
\label{proof21}\s_{12}\mult t_3 =t_3\mult\s_{12}-\s_{32}\mult\s_{13}  {\ds\frac{(\th_{12}+1)}{(\th_{13}+1)\th_{23}}}
+\sum_{k>3}\s_{k2}\mult\s_{1k}X_k^{(3)},
\end{gather}
where all $X_k^{(i)}$ belong to $\Uh$ and the initial $X_k^{(i)}$ are known:
\begin{gather*} X_3^{(1)} =-\frac{\th_{23}+1}{\th_{23}(\th_{13}+1)}  ,\qquad
X_3^{(2)} =\frac{\th_{13}+2}{\th_{23}(\th_{13}+1)}  ,\qquad
X_4^{(3)} =-\frac{\th_{12}+1}{(\th_{14}+1)\th_{24}\th_{34}}  .
\end{gather*}

By the braid group transformation laws, $X_{k+1}^{(i)}=\q_k\big(X_k^{(i)}\big)\cdot(\th_{k,k+1}+1)\th_{k,k+1}^{-1}$
with $k>2$ for $i=1,2$ and $k>3$ for $i=3$, so that
\begin{gather*}%\label{proof22}
X_k^{(1)}=-\frac{1}{\th_{1k}+1}\prod_{j=2}^{k-1}\frac{\th_{jk}+1}{\th_{jk}} ,\qquad
X_k^{(2)}=\frac{\th_{1k}+2}{(\th_{1k}+1)\th_{2k}}\prod_{j=3}^{k-1}
\frac{\th_{jk}+1}{\th_{jk}}  ,\\
X_k^{(3)}=-\frac{\th_{12}+1}{(\th_{1k}+1)\th_{2k}\th_{3k}}\prod_{j=4}^{k-1}
\frac{\th_{jk}+1}{\th_{jk}}  .
\end{gather*}

After the renormalization \rf{not8} and the change of variables~\rf{3.1a}, the relations \rf{proof19}--\rf{proof21} turn into
the relations~\rf{3.12} for $i=1$, $j=2$ and $k=3$.


Applying the transformations from the braid group, we obtain the rest of the relations from the list \rf{3.12}.

{\bf 4a.} We now prove the relations \rf{relation4a} using the arguments similar to \cite[Subsection 6.1.2]{Zh}.
 Consider the
products $H_k\mult_\lambda H_l$  and $H_l\mult_\lambda H_k$ with $\lambda\not=0$.
These products are linear combinations, over $\Uh$, of monomials
\begin{gather*}
a_{kl;\vec{\gamma}}:=H_ke_{-\gamma_1}\cdots e_{-\gamma_m}e_{\gamma_m}\cdots e_{\gamma_1}H_l\qquad\text{and}\qquad
a_{lk;\vec{\gamma}}:=H_le_{-\gamma_1}\cdots e_{-\gamma_m}e_{\gamma_m}\cdots e_{\gamma_1}H_k  ,
\end{gather*}
respectively; here $m\geq 0$ and $\vec{\gamma}:=\{\gamma_1,\dots ,\gamma_m\}$. By construction, the coef\/f\/icient, from $\Uh$, of the monomial
$a_{kl;\vec{\gamma}}$ in $H_k\mult_\lambda H_l$ equals the coef\/f\/icient of $a_{lk;\vec{\gamma}}$ in $H_l\mult_\lambda H_k$.
The expressions $a_{kl;\vec{\gamma}}$ and $a_{lk;\vec{\gamma}}$ are both equal in $\Z_n$  to
\begin{gather*}(\gamma_1,\ve_k)(\gamma_1,\ve_l)E_{-\gamma_1}e_{-\gamma_2}\cdots
e_{-\gamma_m}e_{\gamma_m}\cdots e_{\gamma_2}E_{\gamma_1}  .\end{gather*}
Thus $H_k\mult_\lambda H_l\equiv H_l\mult_\lambda H_k$  for any $\lambda\not=0$. In the zero weight part $\mult_0$
of the product $\mult$ we have the equality $H_kH_l=H_lH_k$ as well. Therefore,
$H_k\mult H_l\equiv H_l\mult H_k$.

{\bf 4b.} The last group \rf{relation4} of relations is left. Like above, we f\/irst explicitly derive the following relation in~$\Z_3$:
\begin{gather} \s_{12}\mult \s_{21}=  h_{12}-t_{12}\mult t_{12}  \frac{1}{\th_{12}-1}+\s_{21}\mult\s_{12}
\frac{(\th_{12}-1)(\th_{12} +2)}{\th_{12}(\th_{12} +1)}\nonumber\\
\phantom{\s_{12}\mult \s_{21}=}{} +  \s_{31}\mult\s_{13} \frac{(\th_{12}-1) (\th_{13} +2)}{\th_{12}\th_{23}(\th_{13}+1)}
-\s_{32}\mult\s_{23}  \frac{\th_{23} +2}{(\th_{23} +1)\th_{13}}  \label{re01}
\end{gather}
(the f\/irst term in the right hand side is $h_{12}$, without hat). The relation \rf{re01} is a corollary of the following calculations
for $\Z_3$, together with the commutation relation $[E_{12},E_{21}]=h_{12}$,
\begin{gather}E_{12}\mult E_{21}\equiv   E_{12}E_{21}-H_{12}^2 \frac{1}{(\th_{12}-1)}
+E_{21}E_{12} \frac{2}{(\th_{12}-1)\th_{12}}\nonumber\\
\phantom{E_{12}\mult E_{21}\equiv}{} - E_{32}E_{23} \frac{\th_{12} -2}{(\th_{12}-1)\th_{13}}+
E_{31}E_{13} \frac{\th_{12} -2}{(\th_{12}-1)\th_{12}\th_{13}} ,\label{zemod1}\\
 H_{12}\mult H_{12}\equiv   H_{12}^2-E_{21}E_{12} \frac{4}{\th_{12} +1}
-E_{32}E_{23} \frac{1}{\th_{23} +1}\nonumber\\
\phantom{H_{12}\mult H_{12}\equiv}{} +  E_{31}E_{13}\left(-1+\frac{1}{\th_{23} +1}+\frac{4}{\th_{12} +1}\right)\frac{1}{\th_{13}+1} ,\label{zemod2}\\
E_{21}\mult E_{12}\equiv E_{21}E_{12}-E_{31}E_{13} \frac{1}{\th_{23}} ,\label{zemod3}\\
 E_{32}\mult E_{23}\equiv  E_{32}E_{23}-E_{31}E_{13} \frac{1}{\th_{12}} ,\label{zemod4}\\
 E_{31}\mult E_{13}\equiv  E_{31}E_{13} .\label{zemod5}
 \end{gather}
Here only the calculation of $E_{12}\mult E_{21}$ deserves a little explanation; by  the analogues of the conditions (i) and (ii),
see paragraph 1 of this subsection, non-trivial contributions to $E_{12}\mult E_{21}$ from the weights 0,
$\ve_1-\ve_2$, $2(\ve_1-\ve_2)$, $\ve_1-\ve_3$ and $2\ve_1-\ve_2-\ve_3$ are possible. By the ABRR equation,
$\J_{2(\ve_1-\ve_2)}(2\th_{12}+4)=-T_{\ve_1-\ve_2}(\J_{\ve_1-\ve_2})$ and $\J_{2\ve_1-\ve_2-\ve_3}
(2\th_1-\th_2-\th_3+3)=-T_{\ve_1-\ve_2}(\J_{\ve_1-\ve_3})-T_{\ve_1-\ve_3}(\J_{\ve_1-\ve_2})-T_{\ve_2-\ve_3}
(\J_{2(\ve_1-\ve_2)})$. We leave further details to the reader.


By the stabilization law in $\Z_4$ we have a relation
\begin{gather} \label{proof23}
\s_{12}\mult \s_{21}= h_{12}-t_{12}\mult t_{12}  \frac{1}{\th_{12}-1}
+\sum_{1\leq i<j\leq n}\s_{ji}\mult\s_{ij}X_{ij} ,\qquad X_{ij}\in\Uh
\end{gather}
with $n=4$, which dif\/fers from \rf{re01} by a presence of terms
\begin{gather*}\s_{43}\mult\s_{34} ,\qquad \s_{42}\mult\s_{24} ,\qquad \s_{41}\mult\s_{14} ,
\end{gather*}
with coef\/f\/icients in $\Uh$. Consider in $\Z_4$ the products  $\s_{12}\mult\s_{21}$, $t_{12}\mult t_{12}$
and $\s_{ji}\mult\s_{ij}$, $1\leq i<j\leq 4$. The weights $(\ve_3-\ve_4)-(\ve_i-\ve_j)$ do not belong
to the cone $\Q_+$ if $1\leq i<j<4$. Thus in the decomposition
\begin{gather*}
E_{ji}\mult E_{ij}\equiv \sum_{k<l}E_{lk}E_{kl}a_{kl} ,\qquad a_{kl}\in\Uh ,\quad 1\leq i<j<4 ,
\end{gather*}
the term with $E_{43}E_{34}$ has a zero coef\/f\/icient, $a_{34}=0$. The same statement holds for the products
$E_{41}\mult E_{14}$ and $E_{42}\mult E_{24}$ since the weights $(\ve_3-\ve_4)-(\ve_i-\ve_4)$ do
not belong to $\Q_+$ for $i=1,2$. Consider the product $E_{12}\mult E_{21}$. Here the term with
$E_{43}E_{34}$ is equal to $E_{12}\mult_{\ve_1-\ve_2+\ve_3-\ve_4}E_{21}$. By \rf{proof3c},
$\J_{\ve_1-\ve_2+\ve_3-\ve_4}=e_{43}e_{21}\otimes e_{12}e_{34}(\th_{12}+1)^{-1}(\th_{34}+1)^{-1}$ and
\begin{gather*}E_{12}\mult_{\ve_1-\ve_2+\ve_3-\ve_4}E_{21}=
E_{12}e_{43}e_{21}e_{12}e_{34}E_{21} \frac{1}{(\th_{12}-1)(\th_{34}+1)}\equiv 0,\end{gather*}
since $[e_{34},E_{21}]=0$ (and $[E_{12},e_{43}]=0$). In the similar manner, the term with
$E_{43}E_{34}$ in $H_{12}\mult H_{12}$ equals $H_{12}\mult_{\ve_3-\ve_4}H_{12}$ and vanishes
since $[e_{34},H_{12}]=0$.

On the other hand, the product $E_{43}\mult E_{34}$
def\/initely contains $E_{43}E_{34}=E_{43}\mult_0 E_{34}$. We thus conclude that the term
$\s_{43}\mult \s_{34}$ is absent in~\rf{proof23}, that is $X_{34}=0$.


For $n>4$, again by the stabilization law, we have a unique relation of the form \rf{proof23}. By uniqueness, it is invariant with respect to the
transformations $\q_3,\q_4,\dots ,\q_{n-1}$ which do not change the product $\s_{12}\mult \s_{21}$. Since $X_{34}=0$, we f\/ind, applying
$\q_4, \q_5, \dots ,\q_{n-1}$,  that  $X_{3j}=0$, $j>3$, wherefrom we further conclude, applying $\q_3,\q_4,\dots ,\q_{j-2}$, that
$X_{ij}=0$, $2<i<j$. We get f\/inally the following relation in $\Z_n$:
\begin{gather} \label{soo1221}
\s_{12}\mult \s_{21}=h_{12}-t_{12}\mult t_{12} \frac{1}{\th_{12}-1}
+\sum_{k=2\lcd n}\s_{k1}\mult\s_{1k}X_{1k}+\sum_{k=3\lcd n}\s_{k2}\mult\s_{2k}X_{2k}
\end{gather}
with known
\begin{gather*}
X_{12}=\frac{(\th_{12}-1)(\th_{12} +2)}{\th_{12}(\th_{12} +1)} ,\qquad
X_{13}=\frac{(\th_{12}-1) (\th_{13} +2)}{\th_{12}\th_{23}(\th_{13}+1)} ,\qquad
X_{23}=-\frac{\th_{23} +2}{(\th_{23} +1)\th_{13}} .
\end{gather*}
Applying to \rf{soo1221} the transformations $\q_3, \q_4,\dots ,\q_{n-1}$ we f\/ind by uniqueness
\begin{gather*}X_{i,k+1}=\frac{\th_{k,k+1}+1}{\th_{k,k+1}} \cdot\q_k(X_{ik}) ,\qquad i=1,2 ;\quad k=3,4,\dots ,n-1 ,
\end{gather*}
and thus
\begin{gather*}
 X_{1k}=\frac{(\th_{12}-1) (\th_{1k} +2)}{\th_{12}\th_{2k}(\th_{1k}+1)}\cdot\prod_{a=3}^{k-1}\frac{\th_{ak}+1}{\th_{ak}}  ,
\qquad X_{2k}=-\frac{\th_{2k} +2}{(\th_{2k} +1)\th_{1k}}\cdot\prod_{a=3}^{k-1}\frac{\th_{ak}+1}{\th_{ak}}
\end{gather*}
for $k>2$.


After the renormalization \rf{not8} and the change of variables~\rf{3.1a}, the relations~\rf{soo1221}  with the
obtained~$X_{1k}$ and $X_{2k}$ turns into the relation~\rf{relation4} for $i=1$ and $j=2$.


Applying the transformations from the braid group, we obtain the rest of the relations from the list~\rf{relation4}.

\subsection{Proof of Theorem \ref{mthe}}\label{proofmain}

For the proof of Theorem \ref{mthe} we just apply the results of~\cite{KO3}, which state that the system ${\mathfrak R}$ is the system of
def\/ining relations once it is satisf\/ied in the algebra $\Z_n$.

\begin{remark}
An attentive look shows that the system $\mathfrak{R}$ is closed under the anti-involution~$\epsilon$; that is, $\epsilon$~transforms any
relation from $\mathfrak{R}$ into a linear over $\Uh$ combination of relations from~$\mathfrak{R}$. Moreover, $\mathfrak{R}$ and
$\epsilon (\mathfrak{R})$ are equivalent over $\Uh$. Indeed, all relations in Section~\rf{sectionDerRel} were derived in
three steps: f\/irst we derive a relation in $\Z_n$ with $n\leq 4$; next by the stabilization principle we extend the derived relation
to~$\Z_n$ with arbitrary~$n$;
and then we f\/ind the whole list of relations of a given (sub)type by applying the braid group transformations (products of the
generators $\q_i$). Due to \rf{invr} one could use $\q_i^{-1}$ instead of $\q_i$ equivalently over
$\Uh$. A~straightforward calculation establishes the equivalence of the extended to arbitrary $n$ lists $\mathfrak{R}$ and
$\epsilon (\mathfrak{R})$ over $\Uh$ for $\Z_n$ with
$n\leq 4$ (this verif\/ication is lengthy for some relations).
Then with the help of~\rf{qeps} we f\/inish the check of the equivalence of $\mathfrak{R}$ and $\epsilon (\mathfrak{R})$ over $\Uh$ for $\Z_n$ with
arbitrary $n$.


Similar arguments establish  the equivalence of $\mathfrak{R}$ and $\omega (\mathfrak{R})$ over $\Uh$; here $\omega$ is the involution def\/ined
in~\rf{not2a}. In~\cite{KO3} this equivalence was obtained dif\/ferently, as a by-product of the equivalence, over $\Uh$, of the system $\mathfrak{R}$ and
the system  \rf{not6} of ordering relations.
\end{remark}
\begin{thebibliography}{99}
\bibitem[1]{ABRR}
Arnaudon D., Buf\/fenoir E., Ragoucy E., Roche Ph.,
Universal solutions of quantum dynamical Yang--Baxter equations,
\href{http://dx.doi.org/10.1023/A:1007498022373}{{\it Lett. Math. Phys.}} {\bf 44} (1998), 201--214,
\href{http://arxiv.org/abs/q-alg/9712037}{q-alg/9712037}.

\bibitem[3]{AST}
Asherova R.M., Smirnov Yu.F., Tolstoy V.N.,
Description of a certain class of projection operators for complex semi-simple Lie algebras,
{\it Mat. Zametki} {\bf 26} (1979), 15--25 (English transl.: \href{http://dx.doi.org/10.1007/BF01140268}{{\it Math. Notes}} {\bf 26} (1979),   499--504).


\bibitem[5]{K}
Khoroshkin S.,
An extremal projector and dynamical twist,
{\it Teoret. Mat. Fiz.} {\bf 139} (2004),  158--176
(English transl.: \href{http://dx.doi.org/10.1023/B:TAMP.0000022749.42512.fd}{{\it Theoret. and Math. Phys.}} {\bf 139} (2004),   582--597).


\bibitem[6]{KO}
 Khoroshkin  S., Ogievetsky O.,
Mickelsson algebras and Zhelobenko operators,
\href{http://dx.doi.org/10.1016/j.jalgebra.2007.04.020}{{\it J.~Algebra}} {\bf 319} (2008), 2113--2165,
\href{http://arxiv.org/abs/math.QA/0606259}{math.QA/0606259}.

\bibitem[7]{KO3}
Khoroshkin  S., Ogievetsky O.,
Diagonal reduction algebras of $\gl$ type,
{\it Funktsional. Anal. i Prilozhen.} {\bf 44} (2010), no.~3, 27--49  (English transl.: \href{http://dx.doi.org/10.1007/s10688-010-0023-0}{{\it Funct. Anal. Appl.}} {\bf 44} (2010),  182--198),
\href{http://arxiv.org/abs/0912.4055}{arXiv:0912.4055}.

\bibitem[10]{Zh}
Zhelobenko D.,
Representations of reductive Lie algebras, Nauka, Moscow, 1994 (in Russian).

\end{thebibliography}
\end{document}
